\documentclass[11pt,letterpaper]{article}
\usepackage{fontspec}
\setmainfont{texgyrepagella}[Extension=.otf,UprightFont=*-regular,BoldFont=*-bold,ItalicFont=*-italic,BoldItalicFont=*-bolditalic]
\usepackage[bidi=default]{babel}
\babelprovide[import,main]{spanish}
\babelprovide[import]{hebrew}
\babelprovide[import]{ancientgreek}
\babelfont[hebrew]{rm}[Path=fonts/,BoldFont=NotoHebrewSerif-700.ttf,Scale=0.95]{NotoHebrewSerif-400.ttf}
\babelfont[ancientgreek]{rm}[Path=fonts/]{GFSPorson.otf}
\newcommand{\heb}[1]{\mbox{\foreignlanguage{hebrew}{#1}}}
\newcommand{\grc}[1]{\foreignlanguage{ancientgreek}{#1}}
\newcommand{\hebverso}[1]{\par\smallskip{\centering\begin{otherlanguage}{hebrew}\large #1\end{otherlanguage}\par}\smallskip}
\newcommand{\SENOR}{\textsc{Señor}}
\usepackage[margin=2.5cm,headheight=14pt]{geometry}
\usepackage{microtype,ragged2e,booktabs,graphicx,xcolor,csquotes,array}
\usepackage[hidelinks]{hyperref}\usepackage{xurl}\urlstyle{same}
\usepackage[most]{tcolorbox}
\usepackage{tikz}\usetikzlibrary{arrows.meta,positioning,calc,decorations.pathreplacing}
\usepackage{titlesec,fancyhdr}
\definecolor{azul}{HTML}{15437C}\definecolor{azuloscuro}{HTML}{0B1F3A}\definecolor{celeste}{HTML}{5694D6}
\definecolor{dorado}{RGB}{196,154,60}\definecolor{papiro}{RGB}{246,239,220}\definecolor{gris}{RGB}{70,78,96}
\titleformat{\section}{\Large\bfseries\color{azul}}{\thesection.}{0.6em}{}
\titleformat{\subsection}{\large\bfseries\color{azuloscuro}}{\thesubsection}{0.6em}{}
\pagestyle{fancy}\fancyhf{}
\fancyhead[L]{\foreignlanguage{spanish}{\small\itshape\color{gris}La sabiduría de Dios}}\fancyhead[R]{\foreignlanguage{spanish}{\small\color{gris}\thepage}}
\renewcommand{\headrulewidth}{0.4pt}
\setlength{\parindent}{1.27cm}\linespread{1.15}\setlength{\emergencystretch}{3em}
\newtcolorbox{cita}[1][]{enhanced,colback=azul!6!white,colframe=azul,boxrule=0pt,leftrule=2.5pt,arc=0pt,
  left=8pt,right=8pt,top=5pt,bottom=5pt,before skip=8pt,after skip=8pt,fontupper=\itshape,#1}
\newcommand{\fuente}[1]{\par\hfill{\upshape\bfseries\small\textcolor{azul}{#1}}}
\newtcolorbox{nota}[1][]{enhanced,colback=dorado!10,colframe=dorado!85!azul,boxrule=0.6pt,arc=2pt,
  fonttitle=\bfseries\small,coltitle=azuloscuro,colbacktitle=dorado!30,fontupper=\small,title={#1}}
\newtcolorbox{filologia}[1][]{enhanced,colback=white,colframe=azul!45,boxrule=0.5pt,arc=2pt,
  fonttitle=\bfseries\small,coltitle=white,colbacktitle=azul!80,fontupper=\small,title={Filología: #1}}
\newcommand{\refapa}[1]{\par{\RaggedRight\hangindent=1.27cm\hangafter=1\noindent #1\par}\vspace{4pt}}
\newcommand{\destacar}[1]{\textbf{\textcolor{azul}{#1}}}
\newcolumntype{L}[1]{>{\RaggedRight\arraybackslash}p{#1}}


\title{\color{azul}\bfseries La multiforme sabiduría de Dios\\[4pt]\large\color{gris}\mdseries ``\SENOR{} con sabiduría fundó la tierra'' (Pr 3:19)}
\author{}\date{}
\begin{document}
\maketitle\thispagestyle{empty}

\begin{abstract}\noindent
La sabiduría de Dios es el primer atributo de la tercera parte de esta obra, la de los atributos que se revelan en la relación de Dios con el ser humano. El capítulo parte del verso que une toda la serie, Jeremías 9:23--24, que comienza negando al sabio la gloria de su sabiduría. Estudia la raíz hebrea \emph{j-k-m}, que no designa ideas abstractas sino destreza práctica (Bezalel, Salomón, Jetro), recorre el espectro de términos que la Reina-Valera 1960 traduce ``sabiduría'' y toma como texto ancla Proverbios 8:22--31, donde la Sabiduría está junto a Dios ``ordenándolo todo'' en la creación. Define la sabiduría de Dios como su capacidad perfecta y activa de obrar con maestría, por la cual diseña, ordena, sostiene y dirige la creación, la providencia y la redención. La sigue en la creación del mundo y del ser humano, en la providencia y en la historia de José, recoge la exposición de Pink sobre los decretos sabios y la mirada judía (\emph{Bereshit Rabá}, Rashi, el Talmud, la liturgia, la literatura del Segundo Templo y Qumrán). Muestra cómo Pablo llama a Cristo ``sabiduría de Dios'', cumple en 1~Corintios 1:31 el verso de Jeremías y describe la sabiduría con palabras de unión: todo coopera, todo subsiste y todo será reunido en Cristo. Termina con una reflexión sobre la sabiduría como sincronía perfecta de la creación y de los atributos de Dios.
\par\smallskip\noindent\textbf{Palabras clave:} sabiduría, \emph{jojmá}, Proverbios 8, \emph{amón}, providencia, José, \emph{sophía}, 1~Corintios 1, Colosenses 1, sincronía.\end{abstract}

\begin{nota}[Nota sobre el texto]
Las citas bíblicas son de la Reina-Valera 1960. Donde esta escribe ``Jehová'', aquí se escribe \SENOR, que traduce el tetragrama \heb{יהוה}, leído en la tradición judía como \emph{Adonay}. El hebreo sigue la BHS; el griego, NA28 y la Septuaginta de Rahlfs-Hanhart. Las fuentes judías (Midrash, Talmud, Misná, Rashi) se citan desde Sefaria y se marcan como precisión. Los libros de la Sabiduría de Salomón y de Ben Sirá se citan en traducción propia desde la Septuaginta, como testigos de la fe judía del Segundo Templo. Cuando la numeración hebrea difiere de la española se da entre corchetes.
\end{nota}

%===================================================================
\section{Introducción: no se alabe el sabio}

El verso que une esta obra comienza hablando de la sabiduría. Antes de decir en qué debe gloriarse el ser humano, Jeremías dice en qué no:

\hebverso{כֹּה אָמַר יְהוָה אַל־יִתְהַלֵּל חָכָם בְּחָכְמָתוֹ\\ הַשְׂכֵּל וְיָדֹעַ אוֹתִי כִּי אֲנִי יְהוָה}
\begin{cita}Así dijo \SENOR: No se alabe el sabio en su sabiduría, ni en su valentía se alabe el valiente, ni el rico se alabe en sus riquezas. Mas alábese en esto el que se hubiere de alabar: en entenderme y conocerme, que yo soy \SENOR, que hago misericordia, juicio y justicia en la tierra; porque estas cosas quiero, dice \SENOR.\fuente{(Jer 9:23--24 [9:22--23])}\end{cita}

La primera gloria que el profeta descarta es la del \emph{jajam} en su \emph{jojmá}, el sabio en su sabiduría. La única gloria legítima es \heb{הַשְׂכֵּל וְיָדֹעַ אוֹתִי} (\emph{haskel veyadoa otí}), ``entender y conocer'' a Dios. El orden no es casual: la sabiduría humana es la primera que debe rendirse, porque es la que con más facilidad se cree autosuficiente. Este capítulo estudia la sabiduría que no se rinde ante nadie, porque no la recibió de nadie: la de Dios.

El libro de Job formula la pregunta de fondo. El capítulo 28 describe al minero que perfora la roca, desvía los ríos y saca a la luz el oro y las piedras preciosas, pero no encuentra la sabiduría:

\begin{cita}Mas ¿dónde se hallará la sabiduría? ¿Dónde está el lugar de la inteligencia?\fuente{(Job 28:12)}\end{cita}

La respuesta llega al final del poema: solo Dios conoce su lugar, porque la usó al crear.

\begin{cita}Dios entiende el camino de ella, y conoce su lugar. Porque él mira hasta los fines de la tierra, y ve cuanto hay bajo los cielos. Al dar peso al viento, y poner las aguas por medida; cuando él dio ley a la lluvia, y camino al relámpago de los truenos, entonces la veía él, y la manifestaba; la preparó y la descubrió también. Y dijo al hombre: He aquí que el temor del Señor es la sabiduría, y el apartarse del mal, la inteligencia.\fuente{(Job 28:23--28)}\end{cita}

Dios ``dio peso al viento'' y puso ``las aguas por medida'': la sabiduría aparece como la destreza con que el Creador calibra el mundo. Al ser humano le toca la sabiduría bajo otra forma, como \emph{temor del Señor}. No la posee como Dios; la recibe en una relación.

Por eso esta serie coloca la sabiduría entre los atributos de la \emph{relación}, los que responden a la pregunta ¿cómo se revela Dios? Como todos los atributos, es una perfección eterna de Dios. Es \emph{incomunicable} en su plenitud, porque solo Dios es ``el único y sabio Dios'' (Ro 16:27) y nadie le enseñó ni lo aconsejó (Is 40:13--14); y es \emph{comunicable} como don, porque ``\SENOR{} da la sabiduría'' (Pr 2:6) y la da ``a todos abundantemente y sin reproche'' (Stg 1:5).

\begin{center}\small\renewcommand{\arraystretch}{1.25}
\begin{tabular}{@{}L{3.4cm}L{5.4cm}L{5.6cm}@{}}
\toprule
\textbf{Mirada} & \textbf{Pregunta} & \textbf{Textos centrales} \\ \midrule
Hebrea & ¿Cómo dice Israel la sabiduría de Dios? & Pr 3:19--20; 8:22--31; Sal 104:24; Jer 10:12; Gn 50:20 \\
Judía & ¿Cómo la leen los sabios de Israel? & \emph{Bereshit Rabá} 1:1; Rashi a Éx 31:3; \emph{Berajot} 58a, 60b; 1QS III \\
Pink & ¿Qué afirma la Escritura y qué exige? & Pink (2020, caps. 2 y 8) \\
Paulina & ¿Qué lugar tiene en Cristo? & 1~Co 1:18--31; 2:7--8; Ro 11:33--36; Ef 1:10; 3:10; Col 1:15--17; 2:3 \\
\bottomrule
\end{tabular}\end{center}

%===================================================================
\section{Texto ancla: Proverbios 8:22--31}

\subsection{La Sabiduría junto al Creador}

En Proverbios 8 la Sabiduría habla en primera persona. Después de llamar a los hombres en las puertas de la ciudad (8:1--21), cuenta su origen: estaba con Dios antes de que existiera el mundo, y estaba presente cuando lo hizo.

\hebverso{יְהוָה קָנָנִי רֵאשִׁית דַּרְכּוֹ קֶדֶם מִפְעָלָיו מֵאָז\\ וָאֶהְיֶה אֶצְלוֹ אָמוֹן וָאֶהְיֶה שַׁעֲשֻׁעִים יוֹם יוֹם}
\begin{cita}\SENOR{} me poseía en el principio, ya de antiguo, antes de sus obras. Eternamente tuve el principado, desde el principio, antes de la tierra. Antes de los abismos fui engendrada; antes que fuesen las fuentes de las muchas aguas. Antes que los montes fuesen formados, antes de los collados, ya había sido yo engendrada; no había aún hecho la tierra, ni los campos, ni el principio del polvo del mundo. Cuando formaba los cielos, allí estaba yo; cuando trazaba el círculo sobre la faz del abismo; cuando afirmaba los cielos arriba, cuando afirmaba las fuentes del abismo; cuando ponía al mar su estatuto, para que las aguas no traspasasen su mandamiento; cuando establecía los fundamentos de la tierra, con él estaba yo ordenándolo todo, y era su delicia de día en día, teniendo solaz delante de él en todo tiempo. Me regocijo en la parte habitable de su tierra; y mis delicias son con los hijos de los hombres.\fuente{(Pr 8:22--31)}\end{cita}

\begin{filologia}[Pr 8:22--31]
\begin{tabular}{@{}L{3.2cm}L{1.6cm}L{10cm}@{}}
\heb{קָנָנִי} \emph{qanani} (v. 22) & \heb{קנה} & ``me adquirió, me poseyó''. La LXX traduce \grc{ἔκτισέν με}, ``me creó'', lectura que pesó en los debates cristológicos de los siglos IV y V. \\[2pt]
\heb{רֵאשִׁית דַּרְכּוֹ} (v. 22) & \heb{ראשׁ} & ``principio de su camino'': la misma palabra con que empieza Gn 1:1 (\heb{בְּרֵאשִׁית}). \\[2pt]
\heb{נִסַּכְתִּי} \emph{nissajti} (v. 23) & \heb{נסך} & ``fui establecida'' (RVR1960: ``tuve el principado''). \\[2pt]
\heb{חוּג} \emph{jug} (v. 27) & \heb{חוג} & ``círculo, bóveda'': Dios traza con compás el horizonte sobre el abismo. \\[2pt]
\heb{חֻקּוֹ} \emph{juqqó} (v. 29) & \heb{חקק} & ``su estatuto'', de ``grabar'': el mar tiene un límite grabado. \\[2pt]
\heb{אָמוֹן} \emph{amón} (v. 30) & \heb{אמן} & ``artesano, maestro de obras'' (cf. \heb{אָמָּן} \emph{ommán}, ``artista'', Cnt 7:1 [7:2]); otros vocalizan \heb{אָמוּן}, ``niño criado''. RVR1960: ``ordenándolo todo''; LXX: \grc{ἁρμόζουσα}, ``ajustando, armonizando''. \\[2pt]
\heb{מְשַׂחֶקֶת} \emph{mesajéqet} (vv. 30--31) & \heb{שׂחק} & ``jugando, regocijándose'': la creación es obra de deleite, no de esfuerzo. \\
\end{tabular}
\end{filologia}

El texto no presenta a una Sabiduría que observa desde afuera. La presenta \emph{junto} a Dios (\heb{אֶצְלוֹ}) mientras traza el horizonte, fija el límite del mar y pone los cimientos de la tierra. El término clave es \emph{amón}. Leído como ``artesano'', describe a la Sabiduría como el maestro de obras del cosmos; leído como ``niño criado'', la describe como la hija en que Dios se deleita. La RVR1960 tradujo el sentido con una frase: ``con él estaba yo ordenándolo todo''. La Septuaginta, siglos antes, había escogido \grc{ἁρμόζουσα}: la que ajusta y pone en armonía las partes. Las dos versiones coinciden en lo esencial: la sabiduría es lo que hace que el mundo no sea un amontonamiento de cosas, sino una obra ordenada.

\begin{nota}[Precisión]
La Sabiduría de Proverbios 8 es una personificación poética de un atributo de Dios. Leerla como figura de Cristo es una lectura tipológica del Nuevo Testamento y de la tradición cristiana (ver sección~\ref{sec:cristo}); el texto de Proverbios no la identifica con una persona divina.
\end{nota}

\subsection{Con sabiduría fundó la tierra}

Proverbios 3 dice lo mismo en forma de sentencia, y Jeremías lo repite casi con las mismas palabras:

\hebverso{יְהוָה בְּחָכְמָה יָסַד־אָרֶץ כּוֹנֵן שָׁמַיִם בִּתְבוּנָה\\ בְּדַעְתּוֹ תְּהוֹמוֹת נִבְקָעוּ וּשְׁחָקִים יִרְעֲפוּ־טָל}
\begin{cita}\SENOR{} con sabiduría fundó la tierra; afirmó los cielos con inteligencia. Con su ciencia los abismos fueron divididos, y destilan rocío los cielos.\fuente{(Pr 3:19--20)}\end{cita}

\hebverso{עֹשֵׂה אֶרֶץ בְּכֹחוֹ מֵכִין תֵּבֵל בְּחָכְמָתוֹ וּבִתְבוּנָתוֹ נָטָה שָׁמָיִם}
\begin{cita}El que hizo la tierra con su poder, el que puso en orden el mundo con su saber, y extendió los cielos con su sabiduría.\fuente{(Jer 10:12)}\end{cita}

\begin{filologia}[Pr 3:19--20 y Jer 10:12]
\begin{tabular}{@{}L{3.2cm}L{1.6cm}L{10cm}@{}}
\heb{בְּחָכְמָה} \emph{bejojmá} & \heb{חכם} & ``con sabiduría'': el instrumento con que Dios ``fundó'' (\heb{יָסַד}) la tierra. \\[2pt]
\heb{בִּתְבוּנָה} \emph{bitvuná} & \heb{בין} & ``con inteligencia, discernimiento'': la capacidad de distinguir y relacionar. \\[2pt]
\heb{בְּדַעְתּוֹ} \emph{beda'tó} & \heb{ידע} & ``con su conocimiento'' dividió los abismos. \\[2pt]
\heb{מֵכִין} \emph{mejín} & \heb{כון} & ``establece, deja a punto'' (RVR1960: ``puso en orden''). En Jer 10:12 la RVR1960 vierte \emph{jojmá} como ``saber'' y \emph{tevuná} como ``sabiduría''. \\
\end{tabular}
\end{filologia}

Sabiduría, inteligencia y conocimiento forman una tríada. La misma tríada vuelve en Éxodo 31:3, cuando Dios llena a Bezalel ``en sabiduría y en inteligencia, en ciencia y en todo arte'' para construir el tabernáculo. La destreza con que Dios hizo el cosmos es la que da al artesano que construye su casa. Jeremías, además, une la sabiduría al poder (\heb{בְּכֹחוֹ}): en Dios la capacidad de hacer y el saber cómo hacerlo no se separan (Job 12:13).

%===================================================================
\section{La palabra}

\subsection{La raíz \heb{ח־כ־ם}}

\emph{Jojmá} (\heb{חָכְמָה}) proviene de la raíz \heb{ח־כ־ם} (\emph{j-k-m}; Strong H2449--H2451). En hebreo bíblico esta raíz no habla de ideas abstractas, sino de habilidad experta, de destreza para hacer algo correctamente, de capacidad de producir un resultado bien hecho. Por eso un artesano, un gobernante o un justo pueden tener \emph{jojmá}. La palabra siempre implica capacidad eficaz con propósito correcto: saber hacer, y hacerlo bien (Koehler y Baumgartner, 2001). En el arameo de Daniel la palabra es la misma (\heb{חָכְמְתָא}, Dn 2:20--23), y la Septuaginta y el Nuevo Testamento la traducen \grc{σοφία} (\emph{sophía}).

\begin{center}\small\renewcommand{\arraystretch}{1.25}
\begin{tabular}{@{}L{2.4cm}L{2.3cm}L{5.6cm}L{4.1cm}@{}}
\toprule
\textbf{Ámbito} & \textbf{Ejemplo} & \textbf{Descripción} & \textbf{Texto} \\ \midrule
Destreza, técnica & Bezalel & hacer correctamente una obra (arte, oficio) bajo dirección divina & Éx 31:3 \\
Gobierno, juicio & Salomón & discernimiento para gobernar y juzgar con justicia & 1~R 3:9, 28 \\
Consejo, dirección & Jetro & ver un problema y dar una solución sabia y ordenada & Éx 18:17--21 \\
Vida justa & el justo (\emph{tsaddiq}) & conducta recta que refleja sabiduría en la vida diaria & Pr 20:7 \\
Hablar sabio & el justo & expresar sabiduría con la boca: enseñanza, consejo & Sal 37:30 \\
Consejo colectivo & los consejeros & seguridad y dirección por muchos consejos & Pr 11:14; 20:18 \\
\bottomrule
\end{tabular}\end{center}

\begin{cita}y lo he llenado del Espíritu de Dios, en sabiduría y en inteligencia, en ciencia y en todo arte, para inventar diseños, para trabajar en oro, en plata y en bronce.\fuente{(Éx 31:3--4)}\end{cita}

\begin{filologia}[Éx 31:3--4, 1~R 3:9 y Sal 37:30]
\begin{tabular}{@{}L{3.2cm}L{1.6cm}L{10cm}@{}}
\heb{לַחְשֹׁב מַחֲשָׁבֹת} \emph{lajshov majashavot} & \heb{חשׁב} & ``para diseñar diseños'' (RVR1960: ``para inventar diseños''). La sabiduría de Bezalel se expresa en planos. \\[2pt]
\heb{לֵב שֹׁמֵעַ} \emph{lev shoméa} & \heb{שׁמע} & ``corazón que escucha'' (RVR1960: ``corazón entendido''): lo que pide Salomón para juzgar (1~R 3:9). \\[2pt]
\heb{יֶהְגֶּה חָכְמָה} \emph{yehgé jojmá} & \heb{הגה} & ``medita, musita sabiduría'': la boca del justo la repite en voz baja (Sal 37:30). \\
\end{tabular}
\end{filologia}

\subsection{``Sabiduría'' en la Reina-Valera 1960}

El material de estudio de esta serie cuenta 226 apariciones de ``sabiduría''. Un conteo propio en la RVR1960, sin los títulos de sección que agregan los editores, da 245 apariciones de ``sabiduría'' o ``sabidurías'' en 233 versículos: 195 en el Antiguo Testamento y 50 en el Nuevo. Si se incluyen los títulos de sección, la cifra sube a 262. La tabla muestra qué palabras del original hay detrás.

\begin{center}\small\renewcommand{\arraystretch}{1.2}
\begin{tabular}{@{}L{5.6cm}L{5.2cm}r@{}}
\toprule
\textbf{Término original} & \textbf{Sentido} & \textbf{Veces} \\ \midrule
\heb{חָכְמָה} \emph{jojmá} (hebreo y arameo) & destreza sabia & 161 \\
\heb{דַּעַת} \emph{da'at}, \heb{דֵּעָה} \emph{de'á} & conocimiento & 24 \\
\heb{שֵׂכֶל} \emph{sékel}, \heb{הִשְׂכִּיל} \emph{hiskil} & perspicacia, éxito & 4 \\
\heb{תּוּשִׁיָּה} \emph{tushiyá} & sabiduría eficaz & 3 \\
\heb{בִּינָה} \emph{biná} & discernimiento & 1 \\
\heb{כֹּחַ לֵב} \emph{kóaj lev} (Job 36:5) & ``fuerza de corazón'' & 1 \\
sin término equivalente (Job 40:2) & --- & 1 \\
\grc{σοφία} \emph{sophía} & sabiduría & 50 \\
\bottomrule
\end{tabular}\end{center}

\begin{nota}[Precisión]
Conteo propio con la RVR1960 y la BHS con números Strong; los casos sin etiqueta Strong se asignaron por la raíz del verso hebreo. Conviene verificarlo en e-Sword: la diferencia con la cifra de 226 puede venir de contar versículos en lugar de apariciones, o de otra edición.
\end{nota}

\subsection{La familia de la sabiduría}

La sabiduría tiene parientes en hebreo. Job 12:13 los reúne: ``Con Dios está la sabiduría y el poder; suyo es el consejo y la inteligencia'' (\heb{עִמּוֹ חָכְמָה וּגְבוּרָה לוֹ עֵצָה וּתְבוּנָה}).

\begin{center}\small\renewcommand{\arraystretch}{1.25}
\begin{tabular}{@{}L{3.0cm}L{1.4cm}L{6.0cm}L{4.0cm}@{}}
\toprule
\textbf{Término} & \textbf{Raíz} & \textbf{Imagen y sentido} & \textbf{Textos} \\ \midrule
\heb{חָכְמָה} \emph{jojmá} & \heb{חכם} & destreza para hacer bien las cosas; el ``arte'' de vivir y de obrar & Pr 3:19; Sal 104:24 \\
\heb{תְּבוּנָה} \emph{tevuná}, \heb{בִּינָה} \emph{biná} & \heb{בין} & de \heb{בֵּין}, ``entre'': distinguir una cosa de otra & Pr 3:19; Is 40:28; 1~R 3:9 \\
\heb{דַּעַת} \emph{da'at} & \heb{ידע} & conocimiento por experiencia y trato personal & Pr 3:20; Sal 19:2 [19:3] \\
\heb{עֵצָה} \emph{etsá} & \heb{יעץ} & consejo, plan, designio & Job 12:13; Is 28:29 \\
\heb{תּוּשִׁיָּה} \emph{tushiyá} & discutida & sabiduría que da resultado; éxito sólido & Is 28:29; Pr 2:7 \\
\heb{מַחֲשָׁבָה} \emph{majashavá} & \heb{חשׁב} & ``diseño, plan'' (Bezalel) y ``pensamiento'' de Dios & Éx 31:4; Is 55:8 \\
\bottomrule
\end{tabular}\end{center}

\begin{cita}También esto salió de \SENOR{} de los ejércitos, para hacer maravilloso el consejo y engrandecer la sabiduría.\fuente{(Is 28:29)}\end{cita}

Isaías 28:23--29 es una parábola del labrador: sabe cuándo arar, qué semilla poner en cada surco y con qué instrumento trillar cada grano. Esa destreza, dice el profeta, ``salió de \SENOR''. La palabra traducida ``sabiduría'' es \heb{תּוּשִׁיָּה} (\emph{tushiyá}): la sabiduría que funciona, que produce el resultado buscado.

\subsection{Del hebreo a Pablo}

La Septuaginta tradujo \emph{jojmá} por \grc{σοφία}: ``Dios con la sabiduría fundó la tierra'' (\grc{ὁ θεὸς τῇ σοφίᾳ ἐθεμελίωσεν τὴν γῆν}, Pr 3:19 LXX). Pablo, formado en las Escrituras de Israel, usa la misma palabra para hablar de Cristo crucificado:

\begin{cita}Porque los judíos piden señales, y los griegos buscan sabiduría; pero nosotros predicamos a Cristo crucificado, para los judíos ciertamente tropezadero, y para los gentiles locura; mas para los llamados, así judíos como griegos, Cristo poder de Dios, y sabiduría de Dios.\fuente{(1~Co 1:22--24)}\end{cita}

La palabra que en Proverbios describe cómo Dios hizo el mundo describe en Pablo cómo Dios lo redime. El puente léxico \heb{חָכְמָה} $\rightarrow$ \grc{σοφία} $\rightarrow$ \grc{θεοῦ σοφία} une las dos obras.

%===================================================================
\section{El Nombre que revela la sabiduría}

En la Escritura el nombre de Dios es la revelación de lo que él es (Éx 33:18--19; 34:5--7). La sabiduría no le da un nombre propio como ``el Santo'', pero la Biblia la une expresamente a su Nombre: Daniel bendice el Nombre porque la sabiduría es suya, Isaías anuncia a un Hijo que se llamará ``Admirable, Consejero'', y Pablo termina Romanos con una doxología ``al único y sabio Dios''.

\subsection{``Sea bendito el nombre de Dios''}
\hebverso{לֶהֱוֵא שְׁמֵהּ דִּי־אֱלָהָא מְבָרַךְ מִן־עָלְמָא וְעַד־עָלְמָא \\ דִּי חָכְמְתָא וּגְבוּרְתָא דִּי לֵהּ־הִיא}
\begin{cita}Y Daniel habló y dijo: Sea bendito el nombre de Dios de siglos en siglos, porque suyos son el poder y la sabiduría. El muda los tiempos y las edades; quita reyes, y pone reyes; da la sabiduría a los sabios, y la ciencia a los entendidos. El revela lo profundo y lo escondido; conoce lo que está en tinieblas, y con él mora la luz.\fuente{(Dn 2:20--22)}\end{cita}

\begin{filologia}[Dn 2:20--22]
\begin{tabular}{@{}L{3.2cm}L{1.6cm}L{10cm}@{}}
\heb{שְׁמֵהּ} \emph{shemeh} & \heb{שׁם} & ``su nombre'' (arameo). Daniel no bendice una idea sobre Dios, sino su Nombre. \\[2pt]
\heb{חָכְמְתָא} \emph{jojmetá} & \heb{חכם} & ``la sabiduría'', forma aramea de \emph{jojmá}. \\[2pt]
\heb{גְּבוּרְתָא} \emph{gevurtá} & \heb{גבר} & ``el poder''. Sabiduría y poder van juntos, como en Job 12:13: ``Con Dios está la sabiduría y el poder''. \\[2pt]
\heb{גָּלֵא} \emph{galé} & \heb{גלה} & ``revela'', ``descubre'' (v.~22): la sabiduría de Dios se conoce cuando él muestra lo que estaba oculto. \\
\end{tabular}
\end{filologia}

El contexto es el sueño de Nabucodonosor. Ningún sabio de Babilonia pudo decir al rey lo que había soñado; Dios lo reveló a Daniel ``en visión de noche'' (Dn 2:19), y Daniel respondió bendiciendo el Nombre. Es la escena de Jeremías 9:23--24 en la corte de Babilonia: los sabios del imperio no pueden alabarse de su sabiduría, y el que sabe algo lo sabe porque Dios se lo ha dado.

\subsection{``Admirable, Consejero''}
\hebverso{וַיִּקְרָא שְׁמוֹ פֶּלֶא יוֹעֵץ אֵל גִּבּוֹר \\ אֲבִיעַד שַׂר־שָׁלוֹם}
\begin{cita}Porque un niño nos es nacido, hijo nos es dado, y el principado sobre su hombro; y se llamará su nombre Admirable, Consejero, Dios Fuerte, Padre Eterno, Príncipe de Paz.\fuente{(Is 9:6 [9:5])}\end{cita}

\begin{filologia}[Is 9:6 (9:5 en hebreo)]
\begin{tabular}{@{}L{3.2cm}L{1.6cm}L{10cm}@{}}
\heb{פֶּלֶא} \emph{pele} & \heb{פלא} & ``maravilla'': lo que sobrepasa la capacidad humana. Es la respuesta del ángel a Manoa: ``¿Por qué preguntas por mi nombre, que es admirable?'' (Jue 13:18). \\[2pt]
\heb{יוֹעֵץ} \emph{yo'ets} & \heb{יעץ} & ``consejero'': el que concibe el plan y lo lleva a término. \\[2pt]
\grc{μεγάλης βουλῆς ἄγγελος} & --- & la Septuaginta traduce ``mensajero del gran consejo'' (Is 9:5 LXX). \\
\end{tabular}
\end{filologia}

Isaías ya había usado las mismas dos raíces para describir al \SENOR{}: ``También esto salió de \SENOR{} de los ejércitos, para hacer maravilloso el consejo (\heb{הִפְלִיא עֵצָה}, \emph{hifli etsá}) y engrandecer la sabiduría'' (Is 28:29). El niño recibe un nombre hecho con el vocabulario de la sabiduría divina.

\begin{nota}[Precisión]
La RVR1960 separa ``Admirable'' y ``Consejero''. Muchos intérpretes leen el hebreo como cuatro pares (\emph{pele-yo'ets, El-gibbor, Avi-ad, Sar-shalom}), es decir, ``Consejero maravilloso''. La relación entre Is 9:6 e Is 28:29 es lectura propia, apoyada en el vocabulario común.
\end{nota}

\subsection{``El Dios de todo saber''}
\hebverso{כִּי אֵל דֵּעוֹת יְהוָה \\ וְלֹא נִתְכְּנוּ עֲלִלוֹת}
\begin{cita}No multipliquéis palabras de grandeza y altanería; Cesen las palabras arrogantes de vuestra boca; Porque el Dios de todo saber es \SENOR{}, Y a él toca el pesar las acciones.\fuente{(1~S 2:3)}\end{cita}

\begin{filologia}[1~S 2:3]
\heb{אֵל דֵּעוֹת} (\emph{El de'ot}) significa literalmente ``Dios de conocimientos''. \emph{De'ot} es el plural de \heb{דֵּעָה}, de la raíz \heb{ידע}, ``conocer''; el plural intensifica (RVR1960: ``Dios de todo saber''). \heb{נִתְכְּנוּ} (\emph{nitkenu}, raíz \heb{תכן}) es ``ser pesadas, medidas'': Dios pesa las acciones. El título lo comparte la sabiduría con la omnisciencia.
\end{filologia}

\begin{nota}[Precisión textual]
En 1~S 2:3 el texto consonántico (\emph{ketiv}) dice \heb{וְלֹא}, ``y no'', y la lectura masorética (\emph{qeré}) indica \heb{וְלוֹ}, ``y a él''. La RVR1960 sigue el \emph{qeré}: ``y a él toca el pesar las acciones''.
\end{nota}

Lo dice Ana en el mismo cántico en que confiesa ``No hay santo como \SENOR{}'' (1~S 2:2). Frente a la arrogancia de quienes hablan con altanería, el Nombre de Dios es ``Dios de conocimientos'': nadie puede engañarlo ni corregirlo.

\subsection{``Al único y sabio Dios'', con Pablo}
\begin{cita}al único y sabio Dios, sea gloria mediante Jesucristo para siempre. Amén.\fuente{(Ro 16:27)}\end{cita}

\begin{filologia}[Ro 16:27]
\grc{μόνῳ σοφῷ θεῷ} (\emph{mónō sophō theō}): ``al solo sabio Dios''. \grc{διὰ Ἰησοῦ Χριστοῦ} (\emph{dia Iēsou Christou}): ``por medio de Jesucristo''. La gloria del Dios sabio pasa por Cristo.
\end{filologia}

Pablo cierra Romanos, la carta que exclamó ``¡Oh profundidad de las riquezas de la sabiduría y de la ciencia de Dios!'' (Ro 11:33), nombrando a Dios por su sabiduría. Y lo hace ``mediante Jesucristo'', a quien llama ``sabiduría de Dios'' (\grc{θεοῦ σοφίαν}, 1~Co 1:24). El Nombre que Daniel bendijo y el Consejero que Isaías anunció se encuentran en él.

\begin{nota}[Precisión textual]
En 1~Ti 1:17 la RVR1960 dice ``al único y sabio Dios'' siguiendo el Texto Recibido; las ediciones críticas (NA27/NA28) leen \grc{μόνῳ θεῷ}, ``al único Dios'', sin ``sabio''. En Ro 16:27, en cambio, ``sabio'' está en todas las ediciones.
\end{nota}

\begin{center}
\small\renewcommand{\arraystretch}{1.2}
\begin{tabular}{@{}L{4.2cm}L{3.6cm}L{3.0cm}L{4.6cm}@{}}
\toprule
\textbf{Nombre o título} & \textbf{Hebreo / griego} & \textbf{Texto} & \textbf{Qué revela} \\ \midrule
``Bendito sea el nombre de Dios'' & \heb{שְׁמֵהּ דִּי־אֱלָהָא} & Dn 2:20 & Suya es la sabiduría \\
``Admirable, Consejero'' & \heb{פֶּלֶא יוֹעֵץ} & Is 9:6 [9:5] & El Hijo con nombre de sabiduría \\
``Maravilloso en el consejo'' & \heb{הִפְלִיא עֵצָה} & Is 28:29 & El \SENOR{} de los ejércitos \\
``Dios de todo saber'' & \heb{אֵל דֵּעוֹת} & 1~S 2:3 & Conoce y pesa las acciones \\
``Único y sabio Dios'' & \grc{μόνῳ σοφῷ θεῷ} & Ro 16:27 & La doxología de Pablo \\
``Sabiduría de Dios'' & \grc{θεοῦ σοφία} & 1~Co 1:24 & Cristo \\ \bottomrule
\end{tabular}\par\smallskip
{\footnotesize\color{gris}\textit{Tabla.} La sabiduría en los nombres de Dios. \textit{Nota.} Elaboración propia.}
\end{center}

%===================================================================
\section{Definición}

\subsection{Qué es la sabiduría según la Escritura}

Tres textos, leídos juntos, muestran tres rostros de la sabiduría. Salomón la ve como \emph{maestría creadora}: Dios fundó la tierra con ella (Pr 3:19) y ella estaba junto a él ``ordenándolo todo'' (Pr 8:30). Daniel la ve como \emph{habilidad para sostener la creación y sus tiempos}:

\hebverso{דִּי חָכְמְתָא וּגְבוּרְתָא דִּי לֵהּ־הִיא\\ וְהוּא מְהַשְׁנֵא עִדָּנַיָּא וְזִמְנַיָּא}
\begin{cita}Sea bendito el nombre de Dios de siglos en siglos, porque suyos son el poder y la sabiduría. El muda los tiempos y las edades; quita reyes, y pone reyes; da la sabiduría a los sabios, y la ciencia a los entendidos.\fuente{(Dn 2:20--21)}\end{cita}

E Isaías la ve como \emph{diseño perfecto} que nadie le enseñó:

\begin{cita}¿Quién enseñó al Espíritu de \SENOR, o le aconsejó enseñándole? ¿A quién pidió consejo para ser avisado? ¿Quién le enseñó el camino del juicio, o le enseñó ciencia, o le mostró la senda de la prudencia?\fuente{(Is 40:13--14)}\end{cita}

La sabiduría, entonces, no es solo conocimiento. Es maestría creadora, habilidad para sostener la creación y conducir sus tiempos, y diseño perfecto de redención y consejo. Aplicada a Dios, la sabiduría es el ``arte'' con el que estructuró el universo; el orden detrás de las leyes de la naturaleza; la coherencia que sostiene y guía la historia; el diseño previo a que existieran el tiempo, el espacio y la materia; y el diseño perfecto de la redención.

\subsection{Definición}

\begin{tcolorbox}[colback=papiro,colframe=azul,boxrule=1pt,arc=3pt]
La \heb{חָכְמָה} (\emph{jojmá}) de Dios es su capacidad perfecta y activa para obrar con maestría, mediante la cual \destacar{diseña, ordena, sostiene y dirige} la creación y la historia conforme a su propósito y para su gloria.
\end{tcolorbox}

En forma desarrollada:

\begin{tcolorbox}[colback=papiro,colframe=azul,boxrule=1pt,arc=3pt]
La \heb{חָכְמָה} (\emph{jojmá}) de Dios es un atributo intrínseco e inseparable de su naturaleza, que consiste en su capacidad perfecta y activa de obrar con maestría, mediante la cual Él ---en la unidad absoluta de su ser--- sincroniza de manera perfecta su omnisciencia y omnipresencia con su poder, amor, misericordia y justicia, haciendo que todos sus atributos operen en perfecta armonía para ordenar, sostener y dirigir la creación, la providencia y la redención, cumpliendo su propósito eterno, manifestando su gloria y obrando el bien óptimo de sus hijos.
\end{tcolorbox}

La fórmula ``en la unidad absoluta de su ser'' es importante. En Dios los atributos no son piezas separadas que alguien deba coordinar desde afuera: cada atributo es Dios mismo, y todos son uno en él. Cuando se dice que la sabiduría ``sincroniza'' los demás atributos, se describe cómo los conocemos nosotros: vemos el poder, la justicia y la misericordia obrar juntos sin contradecirse, y llamamos sabiduría a esa concordia.

\begin{nota}[Precisión]
Las dos definiciones son del autor de la serie. La observación sobre la unidad de los atributos sigue la doctrina clásica de la simplicidad divina (Charnock, 1682/1996; Berkhof, 1938).
\end{nota}

\subsection{Omnisciencia y sabiduría}

La sabiduría se distingue de la omnisciencia, aunque nunca está sin ella.

\begin{center}\small\renewcommand{\arraystretch}{1.3}
\begin{tabular}{@{}L{2.4cm}L{6.0cm}L{6.0cm}@{}}
\toprule
 & \textbf{Omnisciencia} & \textbf{Sabiduría} \\ \midrule
Definición & Dios conoce todas las cosas: pasadas, presentes y futuras. Posee conocimiento perfecto de toda la realidad. & Capacidad divina de usar los mejores medios para lograr los mejores fines, para su gloria. Habilidad perfecta para ordenar todo según su propósito. \\
Pregunta & ¿Qué sabe Dios? & ¿Qué hace Dios con lo que sabe? \\
Imagen & el mapa completo & el camino trazado sobre el mapa \\
Texto & ``conoce lo que está en tinieblas'' (Dn 2:22) & ``Hiciste todas ellas con sabiduría'' (Sal 104:24) \\
\bottomrule
\end{tabular}\end{center}

Pink lo dice al hablar de los decretos de Dios: su sabiduría ``se muestra en la selección de los mejores fines posibles y los medios más adecuados para cumplir sus decretos'' (Pink, 2020, cap. 2). La omnisciencia ve todo; la sabiduría elige y ordena.

%===================================================================
\section{La sabiduría en la creación}

\subsection{``En el principio''}

La primera manifestación de la sabiduría de Dios se observa en la creación misma. La Biblia comienza con una afirmación fundamental:

\hebverso{בְּרֵאשִׁית בָּרָא אֱלֹהִים אֵת הַשָּׁמַיִם וְאֵת הָאָרֶץ}
\begin{cita}En el principio creó Dios los cielos y la tierra.\fuente{(Gn 1:1)}\end{cita}

Desde una perspectiva creacionista, esta frase contiene los elementos fundamentales de la realidad: el \emph{tiempo} (``en el principio''), el \emph{espacio} (``los cielos'') y la \emph{materia} (``la tierra''). El principio marca el inicio de la dimensión temporal, que el ser humano experimenta como pasado, presente y futuro; toda la historia de la creación y de la humanidad ocurre dentro de ella. Los cielos representan la dimensión espacial, con su largo, ancho y alto, que hace posible el movimiento, la relación y la interacción entre las criaturas. La tierra representa la realidad material, que se presenta en estados sólido, líquido y gaseoso. Dios crea esta realidad de manera ordenada, para que exista un universo coherente y funcional donde las criaturas puedan vivir.

\begin{nota}[Precisión]
Leer Gn 1:1 como ``tiempo, espacio y materia'' es una lectura del autor, de raíz creacionista, no un análisis gramatical. En hebreo, ``los cielos y la tierra'' es un merismo: dos extremos que nombran la totalidad, ``todo lo que existe''. Las dos lecturas convergen: nada queda fuera de lo que Dios hizo.
\end{nota}

\subsection{Un sistema ordenado}

\hebverso{מָה־רַבּוּ מַעֲשֶׂיךָ יְהוָה כֻּלָּם בְּחָכְמָה עָשִׂיתָ מָלְאָה הָאָרֶץ קִנְיָנֶךָ}
\begin{cita}¡Cuán innumerables son tus obras, oh \SENOR! Hiciste todas ellas con sabiduría; la tierra está llena de tus beneficios.\fuente{(Sal 104:24)}\end{cita}

\begin{filologia}[Sal 104:24]
\begin{tabular}{@{}L{3.2cm}L{1.6cm}L{10cm}@{}}
\heb{כֻּלָּם בְּחָכְמָה עָשִׂיתָ} & \heb{חכם} & ``todas ellas con sabiduría hiciste'': el adverbio precede al verbo; el acento cae en el \emph{cómo}. \\[2pt]
\heb{קִנְיָנֶךָ} \emph{qinyaneja} & \heb{קנה} & ``tus posesiones, tus criaturas'' (RVR1960: ``tus beneficios''): la misma raíz de Pr 8:22, ``\SENOR{} me \emph{poseía}''. \\
\end{tabular}
\end{filologia}

El Salmo 104 recorre la creación como un sistema: las fuentes riegan los valles, los asnos monteses beben, las aves anidan junto a las aguas, la hierba crece para el ganado, la luna señala los tiempos y el sol conoce su ocaso; el león sale de noche y el hombre trabaja de día (104:10--23). Cada criatura tiene su lugar y su hora. Dios no creó un universo caótico; creó un mundo con estructura, leyes y equilibrio: ciclos de la naturaleza, equilibrio de los ecosistemas, interacción entre las criaturas y dependencia mutua en la vida. Todo eso refleja orden y diseño, y las criaturas fueron creadas dentro de ese sistema para moverse en el tiempo, el espacio y la materia en armonía.

Isaías lleva la imagen al taller del ingeniero:

\begin{cita}¿Quién midió las aguas con el hueco de su mano y los cielos con su palmo, con tres dedos juntó el polvo de la tierra, y pesó los montes con balanza y con pesas los collados?\fuente{(Is 40:12)}\end{cita}

\begin{filologia}[Is 40:12--14, 28]
\begin{tabular}{@{}L{3.2cm}L{1.6cm}L{10cm}@{}}
\heb{מָדַד} \emph{madad} & \heb{מדד} & ``midió'': las medidas de un constructor. \\[2pt]
\heb{תִּכֵּן} \emph{tikkén} & \heb{תכן} & ``ajustó, calibró'' los cielos (v. 12). El mismo verbo abre el v. 13: ``¿quién \emph{calibró} el Espíritu del \SENOR?'' (\heb{מִי־תִכֵּן אֶת־רוּחַ יְהוָה}). \\[2pt]
\heb{שָׁקַל} \emph{shaqal} & \heb{שׁקל} & ``pesó'' con balanza. \\[2pt]
\heb{אֵין חֵקֶר לִתְבוּנָתוֹ} & \heb{חקר} & ``no hay sondeo para su entendimiento'' (v. 28): nadie puede llegar al fondo. \\
\end{tabular}
\end{filologia}

El juego de Isaías es preciso: el Creador calibra el universo, pero nadie calibra al Creador. La sabiduría que dio medida a todo no tiene medida.

\subsection{La creación del ser humano}

La culminación de la creación es el ser humano. La Biblia enseña que Dios no solo lo creó, sino que lo conoce íntimamente desde antes de su nacimiento.

\hebverso{כִּי־אַתָּה קָנִיתָ כִלְיֹתָי תְּסֻכֵּנִי בְּבֶטֶן אִמִּי\\ גָּלְמִי רָאוּ עֵינֶיךָ וְעַל־סִפְרְךָ כֻּלָּם יִכָּתֵבוּ}
\begin{cita}Porque tú formaste mis entrañas; tú me hiciste en el vientre de mi madre. Te alabaré; porque formidables, maravillosas son tus obras; estoy maravillado, y mi alma lo sabe muy bien. No fue encubierto de ti mi cuerpo, bien que en oculto fui formado, y entretejido en lo más profundo de la tierra. Mi embrión vieron tus ojos, y en tu libro estaban escritas todas aquellas cosas que fueron luego formadas, sin faltar una de ellas.\fuente{(Sal 139:13--16)}\end{cita}

\begin{filologia}[Sal 139:13--16]
\begin{tabular}{@{}L{3.2cm}L{1.6cm}L{10cm}@{}}
\heb{קָנִיתָ} \emph{qanita} (v. 13) & \heb{קנה} & ``formaste, adquiriste'': otra vez la raíz de Pr 8:22. \\[2pt]
\heb{תְּסֻכֵּנִי} \emph{tesukkeni} (v. 13) & \heb{סכך} & ``me tejiste'', ``me entrelazaste''. \\[2pt]
\heb{רֻקַּמְתִּי} \emph{ruqqamti} (v. 15) & \heb{רקם} & ``fui bordado'' (RVR1960: ``entretejido''). Es el oficio del bordador del tabernáculo (Éx 35:35). \\[2pt]
\heb{גָּלְמִי} \emph{golmí} (v. 16) & \heb{גלם} & ``mi embrión'': la masa todavía sin forma. \\
\end{tabular}
\end{filologia}

David describe su formación con el vocabulario del taller sagrado: Dios lo tejió y lo bordó como Bezalel y sus artesanos tejían y bordaban las cortinas del tabernáculo. La vida humana no es un accidente biológico, sino parte del diseño consciente de Dios.

La biología describe el mismo proceso desde otra perspectiva. El desarrollo humano comienza en la concepción, cuando una célula del padre y una de la madre se unen y forman una nueva célula, el cigoto, que contiene toda la información genética necesaria para formar un ser humano completo. Después comienza la división celular (mitosis): una célula, dos, cuatro, ocho, dieciséis. A medida que se multiplican, las células se organizan en estructuras cada vez más complejas. Luego ocurre la diferenciación celular: algunas células se convierten en neuronas, otras en células sanguíneas, musculares u óseas; unas forman el cerebro, el corazón, los pulmones, la piel; otras, estructuras pequeñas como uñas, cabello y tejidos de los órganos. Todo el proceso está regulado por la información contenida en el ADN, que dirige el crecimiento, la organización y la función de cada parte del cuerpo.

Lo sorprendente es la coordinación. Las células ``saben'' cuándo dividirse, cuándo detenerse, en qué tejido convertirse y dónde ubicarse dentro del cuerpo. La Biblia interpreta esta realidad como manifestación de la sabiduría divina: Dios diseñó un sistema donde el tiempo permite el desarrollo de la historia, el espacio permite la interacción de la creación y la materia permite la existencia de la vida; y dentro de ese sistema formó al ser humano con una complejidad extraordinaria. ``Mi alma lo sabe muy bien'' (\heb{וְנַפְשִׁי יֹדַעַת מְאֹד}, Sal 139:14).

\begin{nota}[Precisión]
``Saben'' es una manera de hablar: la coordinación celular la regula la información genética y química, no un conocimiento consciente de las células. Que esa regulación exprese la sabiduría de Dios es una lectura de fe, la misma del Salmo 139.
\end{nota}

%===================================================================
\section{La sabiduría en la providencia}

Después de contemplar la sabiduría de Dios en la creación, la Escritura enseña que esa misma sabiduría se manifiesta en la providencia: en la manera en que Dios sostiene, dirige y gobierna la historia y la vida de los seres humanos para cumplir su propósito. La providencia puede definirse como la ejecución del plan de Dios dentro de la historia: la forma en que Dios hace que su voluntad se cumpla en el mundo.

\subsection{Dios gobierna la historia}

\begin{cita}El muda los tiempos y las edades; quita reyes, y pone reyes; da la sabiduría a los sabios, y la ciencia a los entendidos.\fuente{(Dn 2:21)}\end{cita}

Daniel pronuncia estas palabras en Babilonia, en la corte del imperio más poderoso de su tiempo, después de que Dios le revela el sueño de Nabucodonosor. Las palabras arameas \heb{עִדָּנִין} (\emph{iddanín}, ``tiempos'') y \heb{זִמְנִין} (\emph{zimnín}, ``plazos, momentos fijados'') describen a Dios como Señor del calendario de los imperios. La historia no avanza al azar. Dios está por encima de los imperios, de los reyes y de los gobiernos. Las naciones pueden parecer poderosas, pero la autoridad última pertenece a Dios, que decide cuándo se levantan y cuándo terminan los reinados. Los acontecimientos políticos, sociales y culturales de la historia humana están bajo su soberanía providencial.

\subsection{Dios dirige el camino de las personas}

La providencia no se manifiesta solo en la historia mundial, sino también en la vida de cada ser humano.

\hebverso{לֵב אָדָם יְחַשֵּׁב דַּרְכּוֹ וַיהוָה יָכִין צַעֲדוֹ}
\begin{cita}El corazón del hombre piensa su camino; mas \SENOR{} endereza sus pasos.\fuente{(Pr 16:9)}\end{cita}

\begin{filologia}[Pr 16:9]
\begin{tabular}{@{}L{3.2cm}L{1.6cm}L{10cm}@{}}
\heb{יְחַשֵּׁב} \emph{yejashshev} & \heb{חשׁב} & ``piensa, planea'': el hombre diseña su camino, como Bezalel sus planos. \\[2pt]
\heb{יָכִין} \emph{yajín} & \heb{כון} & ``afirma, establece'': el mismo verbo de Jer 10:12 (\heb{מֵכִין תֵּבֵל}). El que ``puso en orden el mundo'' pone en orden los pasos. \\
\end{tabular}
\end{filologia}

El verso describe la experiencia humana: hacemos planes, tomamos decisiones, proyectamos el futuro; sin embargo, Dios dirige finalmente el resultado de nuestras vidas. Esto no significa que el ser humano no tenga responsabilidad, sino que la vida ocurre dentro del marco de la providencia divina.

Pablo lo afirma ante los filósofos de Atenas:

\begin{cita}Y de una sangre ha hecho todo el linaje de los hombres, para que habiten sobre toda la faz de la tierra; y les ha prefijado el orden de los tiempos, y los límites de su habitación.\fuente{(Hch 17:26)}\end{cita}

Dios estableció las épocas de la historia y los lugares donde las personas habitan. No elegimos a nuestros padres, ni el país donde nacemos, ni el tiempo histórico en que vivimos. No decidimos qué acontecimientos globales ocurrirán, y aun en las decisiones políticas o sociales nuestra influencia suele ser mínima. Muchas cosas parecen escapar de nuestro control; en realidad forman parte del gobierno providencial de Dios. Nuestra existencia no es accidental.

\subsection{La voluntad buena, agradable y perfecta}

La providencia tiene un propósito: cumplir la voluntad de Dios.

\begin{cita}No os conforméis a este siglo, sino transformaos por medio de la renovación de vuestro entendimiento, para que comprobéis cuál sea la buena voluntad de Dios, agradable y perfecta.\fuente{(Ro 12:2)}\end{cita}

\begin{center}\small\renewcommand{\arraystretch}{1.25}
\begin{tabular}{@{}L{1.8cm}L{2.6cm}L{6.0cm}L{4.0cm}@{}}
\toprule
\textbf{Español} & \textbf{Griego} & \textbf{Acepciones} & \textbf{Idea central} \\ \midrule
Buena & \grc{ἀγαθός} (\emph{agathós}) & bueno, beneficioso, útil, virtuoso, noble; aquello que produce bienestar real; lo moralmente correcto y provechoso & lo que produce verdadero bien y beneficio para la vida \\
Agradable & \grc{εὐάρεστος} (\emph{euárestos}) & aceptable, grato, que produce satisfacción, que agrada, aprobado por Dios & lo que agrada plenamente a Dios \\
Perfecta & \grc{τέλειος} (\emph{téleios}) & completo, pleno, maduro, acabado; que alcanza su meta (\grc{τέλος}), que cumple su propósito final & lo que alcanza el propósito para el cual fue diseñado \\
\bottomrule
\end{tabular}\end{center}

La providencia de Dios no es arbitraria. Dios dirige la historia y nuestras vidas con sabiduría perfecta, hacia un \emph{télos}. La sabiduría de Dios coordina, sostiene y dirige la historia mediante la providencia para que su voluntad perfecta se cumpla.

\subsection{Dios usa el mal para un bien mayor}

El texto más claro sobre este principio está en Romanos:

\begin{cita}Y sabemos que a los que aman a Dios, todas las cosas les ayudan a bien, esto es, a los que conforme a su propósito son llamados. Porque a los que antes conoció, también los predestinó para que fuesen hechos conformes a la imagen de su Hijo, para que él sea el primogénito entre muchos hermanos.\fuente{(Ro 8:28--29)}\end{cita}

El versículo no dice que todo lo que ocurre es bueno, sino que Dios hace que todo coopere para un bien final dentro de su propósito eterno. Eso incluye el sufrimiento, las pruebas, las injusticias y las dificultades. El verbo es \grc{συνεργεῖ}: ``trabajan juntas'', ``cooperan''. Y el bien final que Dios busca es formar el carácter de sus hijos: transformar su vida, moldear su carácter, hacerlos \grc{συμμόρφους τῆς εἰκόνος τοῦ υἱοῦ αὐτοῦ}, ``conformes a la imagen de su Hijo''.

\begin{nota}[Precisión textual]
Algunos manuscritos antiguos (P\textsuperscript{46}, A, B) leen \grc{πάντα συνεργεῖ ὁ θεός}: ``Dios hace cooperar todas las cosas''. NA28 no incluye \grc{ὁ θεός} en el texto, y la RVR1960 tampoco. Con o sin el sujeto explícito, el contexto (``conforme a su propósito'') atribuye la cooperación a Dios.
\end{nota}

\subsection{El misterio}

La Escritura reconoce también que esta realidad supera nuestra comprensión. La providencia opera en un nivel de sabiduría mucho más amplio que nuestra perspectiva:

\begin{cita}Porque mis pensamientos no son vuestros pensamientos, ni vuestros caminos mis caminos, dijo \SENOR. Como son más altos los cielos que la tierra, así son mis caminos más altos que vuestros caminos, y mis pensamientos más que vuestros pensamientos.\fuente{(Is 55:8--9)}\end{cita}

\begin{cita}Todo lo hizo hermoso en su tiempo; y ha puesto eternidad en el corazón de ellos, sin que alcance el hombre a entender la obra que ha hecho Dios desde el principio hasta el fin.\fuente{(Ec 3:11)}\end{cita}

El Predicador dice \heb{יָפֶה בְעִתּוֹ} (\emph{yafé be-ittó}), ``hermoso en su tiempo''. Lo que hoy parece fuera de lugar tiene su hora en el diseño completo. El ser humano tiene la eternidad en el corazón ---anhela ver el conjunto--- pero no alcanza a ver la obra ``desde el principio hasta el fin''.

%===================================================================
\section{Teología: la exposición de Pink}

A.~W. Pink no dedica un capítulo propio a la sabiduría de Dios; la trata al hablar de los decretos divinos. Su recorrido aporta cuatro ideas. Primera: la Escritura llama ``consejo'' a los decretos de Dios ``para indicar que son perfectamente sabios'', y en los procedimientos divinos ``la sabiduría siempre está asociada con la `voluntad'\,'' (Pink, 2020, cap. 2). En el ser humano la voluntad que no da razones suele ser caprichosa; en Dios, nunca: por eso Pablo habla del ``designio de su voluntad'' (Ef 1:11). Segunda: la sabiduría de Dios ``se muestra en la selección de los mejores fines posibles y los medios más adecuados'' (cap. 2). Tercera: conocemos una parte muy pequeña de sus obras, y debemos ``juzgar el todo por la parte; es decir, lo desconocido por medio de lo conocido''. Pink lo ilustra con una máquina:

\begin{quote}\small
Aquel que percibe el admirable ingenio en el funcionamiento de las partes de una máquina después de haber tenido la oportunidad de examinarla, naturalmente, es conducido a creer que las demás partes de esta máquina son igualmente admirables. (Pink, 2020, cap. 2)
\end{quote}

Cuando lo que no entendemos choca con nuestras nociones de lo bueno y lo sabio, Pink invita a terminar como Pablo: ``¡Oh profundidad de las riquezas de la sabiduría y de la ciencia de Dios!'' (Ro 11:33). Cuarta: los decretos son libres; Dios ``estaba solo cuando hizo sus decretos'' (Is 40:13--14).

En el capítulo sobre la santidad, Pink cita a Stephen Charnock, que une los dos atributos: ``Su justicia es una justicia santa, su sabiduría es una sabiduría santa y su poder es `su santo brazo'\,'' (Pink, 2020, cap. 8). La sabiduría de Dios no es astucia ni mera eficacia. Es eficacia santa: nunca logra un fin bueno por un medio impuro.

\begin{nota}[Precisión]
Pink es aquí una referencia para algunos conceptos; el orden y la estructura del estudio son de la serie.
\end{nota}

%===================================================================
\section{Un encuentro en la historia de Israel: José}

\subsection{Trece años de desorden aparente}

La historia de José expresa la sabiduría de Dios en la providencia con más claridad que ningún otro relato del Antiguo Testamento. Sus hermanos lo venden ``por veinte piezas de plata'' (Gn 37:28); en Egipto es esclavo en casa de Potifar, donde ``\SENOR{} estaba con José'' (39:2); acusado injustamente, termina en la cárcel (39:20); trece años después de ser vendido (37:2; 41:46), interpreta los sueños de Faraón y es puesto ``sobre toda la tierra de Egipto'' (41:41). Ninguno de esos pasos parecía conducir al trono; todos conducían a él. Faraón reconoce de dónde viene su sabiduría:

\hebverso{אַחֲרֵי הוֹדִיעַ אֱלֹהִים אוֹתְךָ אֶת־כָּל־זֹאת אֵין־נָבוֹן וְחָכָם כָּמוֹךָ}
\begin{cita}Y dijo Faraón a José: Pues que Dios te ha hecho saber todo esto, no hay entendido ni sabio como tú.\fuente{(Gn 41:39)}\end{cita}

\subsection{``Dios lo encaminó a bien''}

Cuando el hambre lleva a sus hermanos a Egipto, José lee su propia historia en dos momentos. Al darse a conocer:

\begin{cita}Ahora, pues, no os entristezcáis, ni os pese de haberme vendido acá; porque para preservación de vida me envió Dios delante de vosotros\ldots{} Así, pues, no me enviasteis acá vosotros, sino Dios, que me ha puesto por padre de Faraón y por señor de toda su casa, y por gobernador en toda la tierra de Egipto.\fuente{(Gn 45:5, 8)}\end{cita}

Y después de la muerte de Jacob, cuando sus hermanos temen su venganza:

\hebverso{וְאַתֶּם חֲשַׁבְתֶּם עָלַי רָעָה אֱלֹהִים חֲשָׁבָהּ לְטֹבָה\\ לְמַעַן עֲשֹׂה כַּיּוֹם הַזֶּה לְהַחֲיֹת עַם־רָב}
\begin{cita}Vosotros pensasteis mal contra mí, mas Dios lo encaminó a bien, para hacer lo que vemos hoy, para mantener en vida a mucho pueblo.\fuente{(Gn 50:20)}\end{cita}

\begin{filologia}[Gn 41:39 y 50:20--21]
\begin{tabular}{@{}L{3.2cm}L{1.6cm}L{10cm}@{}}
\heb{נָבוֹן וְחָכָם} \emph{navón vejajam} & \heb{בין}, \heb{חכם} & ``entendido y sabio'': la sabiduría de José es recibida (``Dios te ha hecho saber''). \\[2pt]
\heb{חֲשַׁבְתֶּם} \emph{jashavtem} & \heb{חשׁב} & ``vosotros planeasteis'' el mal. \\[2pt]
\heb{חֲשָׁבָהּ לְטֹבָה} \emph{jashaváh letová} & \heb{חשׁב} & ``Dios lo planeó para bien'': el mismo verbo con otro sujeto. Es el verbo de Bezalel, ``para diseñar diseños'' (Éx 31:4). \\[2pt]
\heb{לְהַחֲיֹת} \emph{lehajayot} & \heb{חיה} & ``para mantener en vida'': el fin del diseño. \\[2pt]
\heb{אֲכַלְכֵּל} \emph{ajalkel} (v. 21) & \heb{כול} & ``yo os sustentaré'': quien vio el diseño puede perdonar y alimentar. \\
\end{tabular}
\end{filologia}

En Génesis 50:20 aparecen dos intenciones simultáneas sobre el mismo hecho: los hermanos actuaron con maldad, y Dios utilizó esa misma situación para salvar a muchas personas. El texto no borra la culpa de los hermanos ni convierte el mal en bien; dice que Dios ``diseñó'' sobre el diseño de ellos. La providencia de Dios puede transformar el mal humano en instrumento de un bien mayor. Ese es el corazón de la sabiduría en la historia: no evita que el hilo se enrede, sino que teje con el hilo enredado.

\begin{nota}[Precisión]
El paralelo entre el \emph{jashav} de Gn 50:20 y el de Bezalel (Éx 31:4) es una observación léxica de este estudio: el verbo es el mismo, pero los textos no se citan entre sí.
\end{nota}

%===================================================================
\section{La sabiduría en el judaísmo}

\subsection{La Torá, plano del Artesano}

El Midrash \emph{Bereshit Rabá} abre su comentario a Génesis con Proverbios 8:30. Explica \heb{אָמוֹן} con cuatro sentidos ---``pedagogo'', ``cubierto'', ``escondido'' y, según algunos, ``grande''--- y luego propone una parábola. Un rey de carne y sangre que construye un palacio no lo hace de su propia idea, sino con un artesano; y el artesano no lo hace de su propia idea, sino que tiene \heb{דִּפְתְּרָאוֹת וּפִנְקְסָאוֹת}, ``rollos y cuadernos'', para saber cómo hacer las habitaciones y las puertas. Así:

\hebverso{הַתּוֹרָה אוֹמֶרֶת אֲנִי הָיִיתִי כְּלִי אֻמְנוּתוֹ שֶׁל הַקָּדוֹשׁ בָּרוּךְ הוּא}

``La Torá dice: yo era el instrumento del oficio del Santo, bendito sea''. El Santo miraba en la Torá y creaba el mundo. Y el Midrash une Gn 1:1 con Pr 8:22: si la Sabiduría es ``principio'' (\heb{רֵאשִׁית}) del camino de Dios, ``en el principio creó Dios'' puede leerse ``con la sabiduría creó Dios'' (\emph{Bereshit Rabá} 1:1).

El Targum Neofiti, una traducción aramea de la Torá usada en la sinagoga, lo dice en el mismo versículo: ``Desde el principio, con sabiduría, el Memra del \SENOR{} creó y perfeccionó los cielos y la tierra'' (Gn 1:1, en McNamara, 1992).

\begin{nota}[Precisión]
Interpretación rabínica. El Midrash identifica la Sabiduría de Proverbios con la Torá, como ya lo hacía Ben Sirá (Sir 24:23). El Nuevo Testamento hará una identificación distinta: con Cristo (sección~\ref{sec:cristo}).
\end{nota}

\subsection{Sabiduría, inteligencia y conocimiento}

Rashi, el gran comentarista medieval, distingue la tríada con que Dios llenó a Bezalel (Éx 31:3):

\begin{center}\small\renewcommand{\arraystretch}{1.3}
\begin{tabular}{@{}L{2.4cm}L{6.2cm}L{5.8cm}@{}}
\toprule
\textbf{Término} & \textbf{Rashi} & \textbf{Traducción} \\ \midrule
\heb{חָכְמָה} & \heb{מַה שֶּׁאָדָם שׁוֹמֵעַ מֵאֲחֵרִים וְלָמֵד} & lo que uno oye de otros y aprende \\
\heb{תְּבוּנָה} & \heb{מֵבִין דָּבָר מִלִּבּוֹ מִתּוֹךְ דְּבָרִים שֶׁלָּמַד} & entiende algo por sí mismo, a partir de lo que aprendió \\
\heb{דַּעַת} & \heb{רוּחַ הַקֹּדֶשׁ} & el Espíritu Santo \\
\bottomrule
\end{tabular}\end{center}

La lectura de Rashi coincide con el texto bíblico en un punto decisivo: la sabiduría más alta del artesano es don del Espíritu (``lo he llenado del Espíritu de Dios'', Éx 31:3).

\subsection{¿Quién es sabio?}

La Misná conserva la definición de Ben Zomá: \heb{אֵיזֶהוּ חָכָם הַלּוֹמֵד מִכָּל אָדָם}, ``¿quién es sabio? El que aprende de todo hombre'', y la apoya en Sal 119:99 (\emph{Pirké Avot} 4:1). El Talmud cuenta que unos sabios del Néguev respondieron a un rey con otra definición: \heb{אֵיזֶהוּ חָכָם הָרוֹאֶה אֶת הַנּוֹלָד}, ``¿quién es sabio? El que ve lo que va a nacer'' (\emph{Tamid} 32a): el que ve el resultado antes de que llegue. En el ser humano es prudencia; en Dios es perfección, porque él anuncia ``lo por venir desde el principio'' (Is 46:10).

\subsection{Bendecir la sabiduría de Dios}

El judaísmo convirtió la sabiduría de Dios en bendición cotidiana. El Talmud enseña (\emph{Berajot} 58a) que quien ve a los sabios de Israel bendice a Dios \heb{שֶׁחָלַק מֵחׇכְמָתוֹ לִירֵאָיו}, ``que repartió de su sabiduría a los que le temen''; y quien ve a los sabios de las naciones, \heb{שֶׁנָּתַן מֵחׇכְמָתוֹ לְבָשָׂר וָדָם}, ``que dio de su sabiduría a carne y sangre''. Toda sabiduría humana, la de Israel y la de los gentiles, es una porción de la sabiduría de Dios. Y quien ve una multitud de Israel bendice al \heb{חֲכַם הָרָזִים}, ``Sabio de los secretos'', porque ``sus mentes no se parecen unas a otras, ni sus rostros'': solo Dios conoce a cada uno.

La bendición más cercana a la reflexión sobre el cuerpo humano es la que se dice cada mañana (\emph{Berajot} 60b):

\hebverso{אֲשֶׁר יָצַר אֶת הָאָדָם בְּחׇכְמָה\\ וּבָרָא בּוֹ נְקָבִים נְקָבִים חֲלָלִים חֲלָלִים}

``\ldots{} que formó al hombre con sabiduría y creó en él orificios y cavidades''. La bendición continúa: si uno de ellos se abriera o se cerrara fuera de su tiempo, sería imposible subsistir y estar delante de Dios. Es el Salmo 139 hecho oración diaria: el cuerpo que funciona a su tiempo es motivo de alabanza cada mañana. Y la oración de la mañana (\emph{Yotser}) recita el Salmo 104:24: ``¡Cuán innumerables son tus obras, \SENOR! Todas las hiciste con sabiduría'' (Elbogen, 1993).

\subsection{La Sabiduría en el judaísmo del Segundo Templo}

Entre el Antiguo y el Nuevo Testamento, la reflexión judía sobre la Sabiduría de Proverbios 8 creció. La \emph{Sabiduría de Salomón}, escrita en griego por un judío de Alejandría (s.~I a.\,e.\,c.), la llama ``artesana de todas las cosas'':

\begin{quote}\small
\grc{ἡ γὰρ πάντων τεχνῖτις ἐδίδαξέ με σοφία}.\\
Porque me enseñó la sabiduría, artesana de todas las cosas. (Sab 7:22)
\end{quote}

\begin{quote}\small
\grc{διατείνει δὲ ἀπὸ πέρατος ἐπὶ πέρας εὐρώστως καὶ διοικεῖ τὰ πάντα χρηστῶς}.\\
Se extiende con fuerza de un extremo al otro, y gobierna todas las cosas con bondad. (Sab 8:1)
\end{quote}

En 7:26 la llama ``resplandor (\grc{ἀπαύγασμα}) de la luz eterna, espejo sin mancha de la actividad de Dios e imagen (\grc{εἰκών}) de su bondad'': un vocabulario que reaparece en Colosenses 1:15 y Hebreos 1:3. Ben Sirá (s.~II a.\,e.\,c.) comienza con ``Toda sabiduría viene del Señor'' (\grc{Πᾶσα σοφία παρὰ κυρίου}, Sir 1:1), y en el capítulo 24 hace que la Sabiduría busque dónde morar hasta plantar su tienda en Jacob; luego la identifica con la Ley: ``Todo esto es el libro de la alianza del Dios Altísimo, la ley que Moisés nos prescribió'' (Sir 24:23).

\begin{nota}[Precisión]
Libros deuterocanónicos, citados en traducción propia desde la LXX (Rahlfs y Hanhart, 2006). Se presentan como testigos de la fe judía del Segundo Templo, no como Escritura de esta serie.
\end{nota}

\subsection{Qumrán: el diseño antes de que existieran}

La \emph{Regla de la Comunidad} de Qumrán abre su tratado de los dos espíritus con una afirmación sobre el plan de Dios:

\hebverso{מֵאֵל הַדֵּעוֹת כּוֹל הוֹוֶה וְנִהְיֶה\\ וְלִפְנֵי הֱיוֹתָם הֵכִין כּוֹל מַחֲשַׁבְתָּם}

``Del Dios de los conocimientos viene todo lo que es y lo que será; y antes de que existieran, estableció todo su diseño; y cuando existen, según su testimonio, cumplen su obra conforme a su glorioso designio, y nada se puede cambiar. En su mano están las leyes de todo, y él los sustenta en todas sus necesidades'' (1QS III 15--17). Las raíces son las de este capítulo: \heb{הֵכִין} (\emph{kun}, ``establecer'', como en Jer 10:12), \heb{מַחֲשָׁבָה} (``diseño'', como en Gn 50:20 y Éx 31:4) y \heb{יְכַלְכְּלֵם} (``los sustenta'', como José en Gn 50:21).

El Gran Rollo de los Salmos de la cueva 11 (11QSal\textsuperscript{a}) contiene un ``Himno al Creador'' desconocido fuera de Qumrán. Sus líneas 13--15 convierten Jeremías 10:12 en bendición: ``Bendito el que hace la tierra con su poder, el que establece el mundo con su sabiduría; con su entendimiento extendió los cielos'' (11QSal\textsuperscript{a} XXVI 13--15).

\begin{nota}[Precisión]
1QS según la transcripción de Qumran-Digital (2022); el Himno al Creador en traducción propia desde la edición de Sanders (1965). Conviene cotejar las imágenes de ambos rollos en la Leon Levy Dead Sea Scrolls Digital Library.
\end{nota}

%===================================================================
\section{La sabiduría en Cristo, leída con Pablo}\label{sec:cristo}

\subsection{Cristo, sabiduría de Dios}

Pablo escribe a una iglesia de Corinto dividida por la admiración a la elocuencia y a la ``sabiduría de palabras'' (1~Co 1:17). Su respuesta es la cruz:

\begin{cita}Pues ya que en la sabiduría de Dios, el mundo no conoció a Dios mediante la sabiduría, agradó a Dios salvar a los creyentes por la locura de la predicación. Porque los judíos piden señales, y los griegos buscan sabiduría; pero nosotros predicamos a Cristo crucificado, para los judíos ciertamente tropezadero, y para los gentiles locura; mas para los llamados, así judíos como griegos, Cristo poder de Dios, y sabiduría de Dios. Porque lo insensato de Dios es más sabio que los hombres, y lo débil de Dios es más fuerte que los hombres.\fuente{(1~Co 1:21--25)}\end{cita}

\begin{filologia}[1~Co 1:21--25]
\begin{tabular}{@{}L{3.2cm}L{1.6cm}L{10cm}@{}}
\grc{μωρία} \emph{mōría} & --- & ``locura, necedad'': así ve el mundo la cruz. \\[2pt]
\grc{σημεῖα\ldots{} σοφίαν} & --- & ``señales\ldots{} sabiduría'': las dos búsquedas, la de Israel y la de Grecia. \\[2pt]
\grc{θεοῦ δύναμιν καὶ θεοῦ σοφίαν} & --- & ``poder de Dios y sabiduría de Dios'': el genitivo antepuesto da énfasis: \emph{de Dios}, no de los hombres. Une los dos atributos que Job 12:13 y Dn 2:20 ponen juntos. \\
\end{tabular}
\end{filologia}

Cristo crucificado responde a las dos búsquedas: es el poder que pedía Israel y la sabiduría que buscaba Grecia, pero no en la forma esperada. La cruz parece locura; es la obra maestra del Artesano, el lugar donde la justicia que condena el pecado y la misericordia que perdona al pecador se encuentran sin anularse (Ro 3:26).

\subsection{El verso de Jeremías cumplido}

Pablo termina el argumento citando el verso que une esta obra:

\begin{cita}Mas por él estáis vosotros en Cristo Jesús, el cual nos ha sido hecho por Dios sabiduría, justificación, santificación y redención; para que, como está escrito: El que se gloría, gloríese en el Señor.\fuente{(1~Co 1:30--31)}\end{cita}

\begin{center}\small\renewcommand{\arraystretch}{1.3}
\begin{tabular}{@{}L{2.4cm}L{6.0cm}L{6.0cm}@{}}
\toprule
\textbf{Texto} & \textbf{Hebreo / griego} & \textbf{Sentido} \\ \midrule
Jer 9:23 & \heb{אַל־יִתְהַלֵּל חָכָם בְּחָכְמָתוֹ} & ``No se alabe el sabio en su sabiduría'' \\
Jer 9:24 LXX (9:23) & \grc{ἐν τούτῳ καυχάσθω ὁ καυχώμενος} & ``en esto gloríese el que se gloría'' \\
1~Co 1:31 & \grc{ὁ καυχώμενος ἐν κυρίῳ καυχάσθω} & resumen paulino de Jeremías \\
\bottomrule
\end{tabular}\end{center}

Jeremías le quitaba al sabio la gloria de su sabiduría; Pablo muestra dónde está la sabiduría que sí permite gloriarse: en Cristo, ``hecho por Dios sabiduría'' para nosotros. Los cuatro sustantivos de 1~Co 1:30 explican en qué consiste esa sabiduría: justificación, santificación y redención. La sabiduría de Dios en Cristo no es una doctrina para admirar, sino una salvación para recibir.

\begin{nota}[Precisión]
La cita de 1~Co 1:31 resume Jer 9:23--24 (LXX 9:22--23); la conexión la reconocen los comentaristas (Hays, 1989).
\end{nota}

\subsection{La sabiduría oculta}

\begin{cita}Mas hablamos sabiduría de Dios en misterio, la sabiduría oculta, la cual Dios predestinó antes de los siglos para nuestra gloria, la que ninguno de los príncipes de este siglo conoció; porque si la hubieran conocido, nunca habrían crucificado al Señor de gloria.\fuente{(1~Co 2:7--8)}\end{cita}

La sabiduría de la cruz fue ``predestinada antes de los siglos'' (\grc{προώρισεν\ldots{} πρὸ τῶν αἰώνων}): el diseño es anterior al tiempo, como la Sabiduría de Proverbios 8:23, establecida ``antes de la tierra''. Los poderes de este siglo creyeron destruir a Cristo y cumplieron el plan: es Génesis 50:20 a escala cósmica. Pablo llama a ese plan \grc{μυστήριον}, no un enigma sin solución, sino un plan escondido y ahora revelado (Ro 16:25--26).

\subsection{``¡Oh profundidad!''}

Romanos 9--11 recorre la historia de Israel: su elección, su tropiezo, el injerto de los gentiles y la promesa de que ``todo Israel será salvo'' (11:26). Pablo no termina con una explicación, sino con adoración:

\begin{cita}¡Oh profundidad de las riquezas de la sabiduría y de la ciencia de Dios! ¡Cuán insondables son sus juicios, e inescrutables sus caminos! Porque ¿quién entendió la mente del Señor? ¿O quién fue su consejero? ¿O quién le dio a él primero, para que le fuese recompensado? Porque de él, y por él, y para él, son todas las cosas. A él sea la gloria por los siglos. Amén.\fuente{(Ro 11:33--36)}\end{cita}

\begin{center}\small\renewcommand{\arraystretch}{1.3}
\begin{tabular}{@{}L{2.0cm}L{12.4cm}@{}}
\toprule
Is 40:13 TM & \heb{מִי־תִכֵּן אֶת־רוּחַ יְהוָה וְאִישׁ עֲצָתוֹ יוֹדִיעֶנּוּ} \\
Is 40:13 LXX & \grc{τίς ἔγνω νοῦν κυρίου, καὶ τίς αὐτοῦ σύμβουλος ἐγένετο} \\
Ro 11:34 & \grc{τίς γὰρ ἔγνω νοῦν κυρίου; ἢ τίς σύμβουλος αὐτοῦ ἐγένετο;} \\
\bottomrule
\end{tabular}\end{center}

\begin{filologia}[Ro 11:33--36]
\begin{tabular}{@{}L{3.2cm}L{1.6cm}L{10cm}@{}}
\grc{βάθος} \emph{báthos} & --- & ``profundidad'', como un abismo sin fondo. \\[2pt]
\grc{ἀνεξεραύνητα} \emph{anexeraúnēta} & --- & ``imposibles de rastrear'': juicios que no se pueden investigar hasta el fondo. \\[2pt]
\grc{ἀνεξιχνίαστοι} \emph{anexijníastoi} & --- & ``sin huellas'': caminos de los que no se pueden seguir las pisadas. \\[2pt]
\grc{ἐξ\ldots{} δι᾽\ldots{} εἰς} & --- & ``de él, por él, para él'': origen, medio y meta de todas las cosas. \\
\end{tabular}
\end{filologia}

Pablo cita Isaías 40:13 según la Septuaginta y recuerda Job 41:11 [41:3] (\heb{מִי הִקְדִּימַנִי וַאֲשַׁלֵּם}, ``¿quién me ha dado a mí primero, para que yo restituya?''). El Dios que nadie calibró (Is 40:12--13) es el que conduce la historia de Israel y de las naciones hacia su fin. La reacción correcta ante su sabiduría no es la explicación exhaustiva, sino la doxología.

\subsection{La multiforme sabiduría de Dios}

\begin{cita}y de aclarar a todos cuál sea la dispensación del misterio escondido desde los siglos en Dios, que creó todas las cosas; para que la multiforme sabiduría de Dios sea ahora dada a conocer por medio de la iglesia a los principados y potestades en los lugares celestiales, conforme al propósito eterno que hizo en Cristo Jesús nuestro Señor.\fuente{(Ef 3:9--11)}\end{cita}

\begin{filologia}[Ef 3:10]
\begin{tabular}{@{}L{3.2cm}L{1.6cm}L{10cm}@{}}
\grc{πολυποίκιλος} \emph{polypoíkilos} & --- & ``de muchos colores, ricamente variada''. \grc{ποικίλος} describe tejidos bordados de varios colores; en la LXX, la túnica de José es \grc{χιτὼν ποικίλος} (Gn 37:3). \\[2pt]
\grc{οἰκονομία} \emph{oikonomía} & --- & ``administración, plan de una casa'' (RVR1960: ``dispensación''). \\[2pt]
\grc{πρόθεσις τῶν αἰώνων} & --- & ``propósito de los siglos'': el designio eterno. \\
\end{tabular}
\end{filologia}

El ``misterio'' que Pablo explica en Efesios 3:6 es que los gentiles son ``coherederos y miembros del mismo cuerpo'' con Israel. La iglesia, hecha de judíos y gentiles reconciliados, es la pieza bordada de muchos colores que exhibe la sabiduría de Dios ante los poderes celestiales. La sabiduría no se exhibe en una idea, sino en un pueblo.

\subsection{Todo reunido en Cristo}

En Efesios 1 y Colosenses 1, Pablo describe la sabiduría de Dios en Cristo con palabras de unión:

\begin{cita}que hizo sobreabundar para con nosotros en toda sabiduría e inteligencia, dándonos a conocer el misterio de su voluntad, según su beneplácito, el cual se había propuesto en sí mismo, de reunir todas las cosas en Cristo, en la dispensación del cumplimiento de los tiempos, así las que están en los cielos, como las que están en la tierra.\fuente{(Ef 1:8--10)}\end{cita}

\begin{cita}El es la imagen del Dios invisible, el primogénito de toda creación. Porque en él fueron creadas todas las cosas, las que hay en los cielos y las que hay en la tierra, visibles e invisibles\ldots{} todo fue creado por medio de él y para él. Y él es antes de todas las cosas, y todas las cosas en él subsisten.\fuente{(Col 1:15--17)}\end{cita}

\begin{cita}en quien están escondidos todos los tesoros de la sabiduría y del conocimiento.\fuente{(Col 2:3)}\end{cita}

\begin{filologia}[Ro 8:28, Col 1:17 y Ef 1:10]
\begin{tabular}{@{}L{3.2cm}L{1.6cm}L{10cm}@{}}
\grc{συνεργεῖ} \emph{synergeî} & \grc{σύν} & ``trabajan juntas, cooperan'' (Ro 8:28): todas las cosas en la historia. \\[2pt]
\grc{συνέστηκεν} \emph{synéstēken} & \grc{σύν} & ``se mantienen unidas, subsisten juntas'' (Col 1:17): todas las cosas en la creación. En la LXX de Job 28:23, Dios ``bien compuso'' (\grc{εὖ συνέστησεν}) el camino de la sabiduría. \\[2pt]
\grc{ἀνακεφαλαιώσασθαι} \emph{anakephalaiōsasthai} & \grc{κεφαλή} & ``reunir bajo una cabeza, recapitular'' (Ef 1:10): todas las cosas al final de los tiempos. \\
\end{tabular}
\end{filologia}

Lo que coopera en la historia (Ro 8:28), lo que se mantiene unido en la creación (Col 1:17) y lo que será reunido al final (Ef 1:10) tienen el mismo centro: Cristo. Pablo atribuye a Cristo lo que Israel decía de la Sabiduría: estaba antes de todas las cosas (Pr 8:22--23), por ella fueron hechas (Pr 3:19), es imagen de la bondad de Dios (Sab 7:26). Por eso puede decir que en él ``están escondidos todos los tesoros de la sabiduría'' (Col 2:3), y que en él ``habita corporalmente toda la plenitud de la Deidad'' (Col 2:9): la tesis de esta obra.

\begin{nota}[Precisión]
La lectura de Col 1:15--20 como cristología de la Sabiduría es amplia en la investigación (Dunn, 1998; Wright, 2013). La autoría paulina de Efesios y Colosenses se discute en la academia; la misma enseñanza aparece en cartas indiscutidas (Ro 11:33--36; 1~Co 1--2; 8:6).
\end{nota}

%===================================================================
\section{Aplicación: vivir a la luz de la sabiduría de Dios}

\subsection{Cómo se adquiere la sabiduría}

La sabiduría de Dios es incomunicable en su plenitud, pero Dios la comparte. La Escritura señala varios caminos.

\begin{enumerate}
\item \textbf{El temor del Señor.} ``El temor de \SENOR{} es el principio de la sabiduría, y el conocimiento del Santísimo es la inteligencia'' (Pr 9:10). La sabiduría empieza por reconocer quién es Dios: no es una técnica, sino una relación (Job 28:28).
\item \textbf{Pedirla.} ``Y si alguno de vosotros tiene falta de sabiduría, pídala a Dios, el cual da a todos abundantemente y sin reproche, y le será dada'' (Stg 1:5). Salomón la pidió y la recibió (1~R 3:9--12).
\item \textbf{La Palabra.} ``El testimonio de \SENOR{} es fiel, que hace sabio al sencillo'' (Sal 19:7 [19:8]); ``me has hecho más sabio que mis enemigos con tus mandamientos'' (Sal 119:98); ``la exposición de tus palabras alumbra; hace entender a los simples'' (Sal 119:130); las Escrituras ``te pueden hacer sabio para la salvación por la fe que es en Cristo Jesús'' (2~Ti 3:15). El \heb{פֶּתִי} (\emph{petí}) de Sal 19:8 es el ingenuo que se deja llevar por cualquiera: la Palabra le da criterio.
\item \textbf{Contar los días.} ``Enséñanos de tal modo a contar nuestros días, que traigamos al corazón sabiduría'' (Sal 90:12). La sabiduría crece al ver la vida desde su final.
\end{enumerate}

\subsection{Dos sabidurías}

Santiago distingue la sabiduría que viene de abajo y la que viene de arriba:

\begin{cita}porque esta sabiduría no es la que desciende de lo alto, sino terrenal, animal, diabólica.\fuente{(Stg 3:15)}\end{cita}

\begin{cita}Pero la sabiduría que es de lo alto es primeramente pura, después pacífica, amable, benigna, llena de misericordia y de buenos frutos, sin incertidumbre ni hipocresía.\fuente{(Stg 3:17)}\end{cita}

La sabiduría ``de lo alto'' se reconoce por su carácter, no por su brillo: es pura antes que eficaz, como la ``sabiduría santa'' de Charnock. La otra produce ``celos y contención'' (Stg 3:16), y es la que Pablo describe en Romanos 1:22: ``Profesando ser sabios, se hicieron necios''.

\subsection{Quién la da a conocer}

Dios encargó a su pueblo difundir su sabiduría. La \emph{iglesia} la da a conocer incluso a los poderes celestiales (Ef 3:10). La \emph{enseñanza mutua} la transmite: ``La palabra de Cristo more en abundancia en vosotros, enseñándoos y exhortándoos unos a otros en toda sabiduría'' (Col 3:16). La \emph{predicación} la anuncia: ``a quien anunciamos, amonestando a todo hombre, y enseñando a todo hombre en toda sabiduría'' (Col 1:28). Y la \emph{oración} la pide para otros: ``os dé espíritu de sabiduría y de revelación en el conocimiento de él'' (Ef 1:17). Al final de la historia, la sabiduría vuelve a quien le pertenece: ``El Cordero que fue inmolado es digno de tomar el poder, las riquezas, la sabiduría, la fortaleza, la honra, la gloria y la alabanza'' (Ap 5:12).

\subsection{Cuando duele}

La sabiduría de Dios no siempre se ve \emph{en} el momento; muchas veces se ve \emph{desde} el final. José esperó trece años para decir ``Dios lo encaminó a bien''. Pablo escribió Romanos 8:28 a creyentes que sufrían ``tribulación, angustia, persecución'' (Ro 8:35). El creyente que no entiende lo que le ocurre puede confiar, porque conoce a quien entiende: ``Dios entiende el camino de ella'' (Job 28:23).

\begin{cita}Fíate de \SENOR{} de todo tu corazón, y no te apoyes en tu propia prudencia. Reconócelo en todos tus caminos, y él enderezará tus veredas.\fuente{(Pr 3:5--6)}\end{cita}

``No te apoyes'' en tu \heb{בִּינָה} (\emph{biná}): la inteligencia humana es buena, pero no es un fundamento. El fundamento es el Dios sabio.

%===================================================================
\section{Reflexión final: la sincronía perfecta}

\subsection{La naturaleza que se orquesta}

Al atardecer, miles de estorninos giran sobre un río como un solo cuerpo. Se estiran, se contraen, cambian de dirección en un instante, y ninguno choca con otro. Nadie los dirige. Un estudio de campo en Roma mostró que cada ave no atiende a toda la bandada, sino a seis o siete vecinos, y que esa regla sencilla basta para que miles se muevan juntas (Ballerini et al., 2008). Las abejas van y vienen de la colmena, y la danza de una obrera indica a las demás la dirección y la distancia del alimento (von Frisch, 1967). Los salmones remontan el río en que nacieron, guiados por el olor de sus aguas (Hasler y Scholz, 1983). En algunos bosques, miles de luciérnagas encienden su luz al mismo tiempo, sin director (Strogatz, 2003). El río baja al mar, el agua sube en vapor, cae en lluvia y vuelve al río.

\begin{cita}Los ríos todos van al mar, y el mar no se llena; al lugar de donde los ríos vinieron, allí vuelven para correr de nuevo.\fuente{(Ec 1:7)}\end{cita}

La naturaleza se orquesta: cada movimiento ocurre en su tiempo, uno al lado del otro, sin hacerse daño. Es una sincronía perfecta. La Escritura ya lo había observado. De las langostas dice que ``no tienen rey, y salen todas por cuadrillas'' (Pr 30:27), y la Septuaginta precisa que marchan ``a una sola orden, en buen orden'' (\grc{ἀφ᾽ ἑνὸς κελεύσματος εὐτάκτως}). De las aves migratorias dice:

\hebverso{גַּם־חֲסִידָה בַשָּׁמַיִם יָדְעָה מוֹעֲדֶיהָ}
\begin{cita}Aun la cigüeña en el cielo conoce su tiempo, y la tórtola y la grulla y la golondrina guardan el tiempo de su venida; pero mi pueblo no conoce el juicio de \SENOR.\fuente{(Jer 8:7)}\end{cita}

Y de las criaturas más pequeñas:

\hebverso{אַרְבָּעָה הֵם קְטַנֵּי־אָרֶץ וְהֵמָּה חֲכָמִים מְחֻכָּמִים}
\begin{cita}Cuatro cosas son de las más pequeñas de la tierra, y las mismas son más sabias que los sabios.\fuente{(Pr 30:24)}\end{cita}

\begin{filologia}[Jer 8:7 y Pr 30:24]
\begin{tabular}{@{}L{3.2cm}L{1.6cm}L{10cm}@{}}
\heb{חֲסִידָה} \emph{jasidá} & \heb{חסד} & ``cigüeña'': lleva la raíz de \emph{jésed}, ``la leal''. \\[2pt]
\heb{מוֹעֲדֶיהָ} \emph{moadeha} & \heb{יעד} & ``sus tiempos señalados'': la palabra de las fiestas de Israel. Las aves guardan sus tiempos; el pueblo olvida los de Dios. \\[2pt]
\heb{חֲכָמִים מְחֻכָּמִים} & \heb{חכם} & ``sabias, hechas sabias'': el participio pasivo (\emph{pual}) dice que su sabiduría es recibida. \\
\end{tabular}
\end{filologia}

Para ver esta sincronía, el video \emph{Murmuration} de Sophie Windsor Clive y Liberty Smith (2011) muestra una bandada de estorninos sobre el río Shannon, en Irlanda, filmada desde una canoa (\url{https://vimeo.com/31158841}).

\subsection{Más que la naturaleza}

Pero la sabiduría de Dios es más que la naturaleza. ``Los cielos cuentan la gloria de Dios, y el firmamento anuncia la obra de sus manos'' (Sal 19:1 [19:2]); la creación \emph{cuenta} la gloria. La sabiduría, sin embargo, se ve todavía más en la creación del ser humano, tejido y bordado en el vientre de su madre (Sal 139); en la soberanía de Dios, que ``quita reyes, y pone reyes'' (Dn 2:21); en la manera en que Dios permite o usa lo que está mal para transformarlo en bien, como en la historia de José (Gn 50:20); y, sobre todo, en el plan de redención: Cristo crucificado, ``sabiduría de Dios'' (1~Co 1:24), y una iglesia de judíos y gentiles que exhibe su ``multiforme sabiduría'' (Ef 3:10).

\subsection{El sincronizador de los atributos}

La sabiduría es superlativa. Es el sincronizador perfecto de todos los atributos de Dios: por ella el poder nunca actúa sin amor, la justicia nunca sin misericordia, la misericordia nunca contra la santidad, la fidelidad nunca sin verdad. En la cruz se ve con más claridad que en cualquier otro lugar: Dios es ``justo, y el que justifica'' (Ro 3:26). Junto con la santidad, la sabiduría es un atributo que se ve en todos los demás. Si la santidad es su \emph{brillo} ---``la pureza es el esplendor de cada atributo'', dice Charnock (en Pink, 2020, cap. 8)---, la sabiduría es su \emph{concierto}.

\begin{nota}[Precisión]
``Sincronizador de los atributos'' es una lectura propia de la serie. Es una analogía: en Dios no hay piezas que deban ajustarse, porque cada atributo es Dios mismo (Charnock, 1682/1996). La imagen describe cómo percibimos la unidad de Dios cuando lo vemos obrar.
\end{nota}

\subsection{Una definición gráfica}

\begin{tcolorbox}[colback=papiro,colframe=azul,boxrule=1pt,arc=3pt]
La sabiduría de Dios es la perfección eterna por la cual Él \destacar{concibe}, \destacar{ordena} y \destacar{conduce} todas las cosas ---desde el átomo hasta la historia de la redención--- hacia su fin, de modo que cada parte se mueve en su tiempo y en su lugar, en concierto con todas las demás, sin choque ni desperdicio; y todos sus atributos obran en Él como una sola voluntad: reunir todas las cosas en Cristo para su gloria (Pr 3:19; Ro 11:36; Ef 1:10; Col 1:17).
\end{tcolorbox}

Si el universo fuera una orquesta, la sabiduría de Dios no sería un instrumento más: sería a la vez la \emph{partitura}, el diseño eterno escrito antes de que existieran el tiempo, el espacio y la materia (Pr 8:22--23; 1~Co 2:7); la \emph{batuta}, la providencia que marca a cada criatura su entrada y su tiempo (Dn 2:21; Ro 8:28); y el \emph{acorde final}, la redención, cuando todas las cosas sean reunidas bajo una sola cabeza, Cristo (Ef 1:10). ``Sincronía'' viene del griego \grc{σύν} y \grc{χρόνος}: tiempo compartido. Eso es lo que hace la sabiduría de Dios con todo lo que existe.

%===================================================================
\section{Síntesis}

\begin{tcolorbox}[colback=papiro,colframe=azul,boxrule=1pt,arc=3pt]
La \emph{jojmá} de Dios es destreza perfecta: con ella fundó la tierra (Pr 3:19), estuvo ``ordenándolo todo'' (Pr 8:30), formó al ser humano (Sal 139), gobierna los tiempos (Dn 2:21) y convirtió en bien el mal hecho a José (Gn 50:20). Israel la bendice cada día; Pablo la ve en Cristo crucificado, ``poder de Dios y sabiduría de Dios'' (1~Co 1:24), en quien todas las cosas subsisten (Col 1:17) y serán reunidas (Ef 1:10). Quien la recibe deja de gloriarse en su propia sabiduría y se gloría en el Señor (Jer 9:23--24; 1~Co 1:31).
\end{tcolorbox}

\begin{center}\small\renewcommand{\arraystretch}{1.3}
\begin{tabular}{@{}L{3.0cm}L{11.4cm}@{}}
\toprule
\textbf{Mirada} & \textbf{Aporte al atributo} \\ \midrule
Hebrea & \emph{Jojmá} como destreza (Bezalel, Salomón, Jetro); Pr 8 y el \emph{amón} que lo ordena todo; Pr 3:19--20 y Jer 10:12; Sal 104 y 139; José y el verbo \emph{jashav}. \\
Judía & La Torá como plano del Artesano (\emph{Bereshit Rabá} 1:1; Targum Neofiti); Rashi; ``el que ve lo que va a nacer''; las bendiciones por la sabiduría y \emph{Asher yatsar}; Sabiduría de Salomón y Ben Sirá; 1QS y el Himno al Creador. \\
Pink & Los decretos son sabios: los mejores fines y los medios más adecuados; la máquina; la sabiduría santa. \\
Paulina & Cristo, poder y sabiduría de Dios; Jer 9:24 en 1~Co 1:31; la sabiduría oculta; Ro 11:33--36; la multiforme sabiduría (Ef 3:10); todo coopera, subsiste y se reúne en Cristo. \\
\bottomrule
\end{tabular}\end{center}

\begin{cita}Mas alábese en esto el que se hubiere de alabar: en entenderme y conocerme, que yo soy \SENOR, que hago misericordia, juicio y justicia en la tierra; porque estas cosas quiero, dice \SENOR.\fuente{(Jer 9:24 [9:23])}\end{cita}

El verso que une esta obra empieza quitándole la gloria al sabio y termina dándosela al que conoce a Dios. Conocer a Dios es conocer al único sabio.

%===================================================================
\section*{Glosario}

\begin{description}\setlength{\itemsep}{2pt}
\item[\emph{Amón}] (\heb{אָמוֹן}) artesano, maestro de obras; o niño criado (Pr 8:30).
\item[\emph{Anakephalaiōsis}] (\grc{ἀνακεφαλαίωσις}) recapitulación: reunir todo bajo una cabeza (Ef 1:10).
\item[\emph{Biná}] (\heb{בִּינָה}) discernimiento, entendimiento.
\item[\emph{Da'at}] (\heb{דַּעַת}) conocimiento por experiencia.
\item[\emph{Etsá}] (\heb{עֵצָה}) consejo, designio.
\item[\emph{Jajam}] (\heb{חָכָם}) sabio, hábil.
\item[\emph{Jashav}] (\heb{חָשַׁב}) pensar, planear, diseñar (Gn 50:20; Éx 31:4).
\item[\emph{Jojmá}] (\heb{חָכְמָה}) sabiduría, destreza.
\item[\emph{Majashavá}] (\heb{מַחֲשָׁבָה}) diseño, plan, pensamiento.
\item[\emph{Mystērion}] (\grc{μυστήριον}) plan oculto y ahora revelado (1~Co 2:7).
\item[\emph{Polypoíkilos}] (\grc{πολυποίκιλος}) multiforme, de muchos colores (Ef 3:10).
\item[\emph{Sophía}] (\grc{σοφία}) sabiduría.
\item[\emph{Synergéō}] (\grc{συνεργέω}) cooperar (Ro 8:28).
\item[\emph{Synístēmi}] (\grc{συνίστημι}) mantener unido, subsistir (Col 1:17).
\item[\emph{Tevuná}] (\heb{תְּבוּנָה}) inteligencia, entendimiento.
\item[\emph{Tushiyá}] (\heb{תּוּשִׁיָּה}) sabiduría eficaz (Is 28:29).
\end{description}

%===================================================================
\section*{Referencias}

\refapa{Ballerini, M., Cabibbo, N., Candelier, R., Cavagna, A., Cisbani, E., Giardina, I., Lecomte, V., Orlandi, A., Parisi, G., Procaccini, A., Viale, M., \& Zdravkovic, V. (2008). Interaction ruling animal collective behavior depends on topological rather than metric distance: Evidence from a field study. \emph{Proceedings of the National Academy of Sciences, 105}(4), 1232--1237. \url{https://doi.org/10.1073/pnas.0711437105}}
\refapa{Bauer, W., Danker, F.~W., Arndt, W.~F., \& Gingrich, F.~W. (2000). \emph{A Greek-English lexicon of the New Testament and other early Christian literature} (3.ª ed.). University of Chicago Press.}
\refapa{Berkhof, L. (1938). \emph{Systematic theology}. Eerdmans.}
\refapa{Charnock, S. (1996). \emph{The existence and attributes of God} (Vols. 1--2). Baker Books. (Obra original publicada en 1682)}
\refapa{Dunn, J.~D.~G. (1998). \emph{The theology of Paul the apostle}. Eerdmans.}
\refapa{Elbogen, I. (1993). \emph{Jewish liturgy: A comprehensive history} (R.~P. Scheindlin, Trad.). Jewish Publication Society.}
\refapa{Elliger, K., \& Rudolph, W. (Eds.). (1997). \emph{Biblia Hebraica Stuttgartensia} (5.ª ed.). Deutsche Bibelgesellschaft.}
\refapa{Hasler, A.~D., \& Scholz, A.~T. (1983). \emph{Olfactory imprinting and homing in salmon}. Springer.}
\refapa{Hays, R.~B. (1989). \emph{Echoes of Scripture in the letters of Paul}. Yale University Press.}
\refapa{Koehler, L., \& Baumgartner, W. (2001). \emph{The Hebrew and Aramaic lexicon of the Old Testament} (M.~E.~J. Richardson, Trad.; Ed. de estudio). Brill.}
\refapa{McNamara, M. (1992). \emph{Targum Neofiti 1: Genesis} (The Aramaic Bible, Vol. 1A). Liturgical Press.}
\refapa{\emph{Misná: Pirké Avot} [4:1]. (s.f.). Sefaria. \url{https://www.sefaria.org}}
\refapa{Nestle, E., \& Aland, K. (Eds.). (2012). \emph{Novum Testamentum Graece} (28.ª ed.). Deutsche Bibelgesellschaft.}
\refapa{Pink, A.~W. (2020). \emph{Los atributos de Dios} (E. Rodríguez Cruz, Trad.). Chapel Library. (Obra original publicada en 1930)}
\refapa{Qumran-Digital. (2022). \emph{1QS (transcripción)}. Akademie der Wissenschaften zu Göttingen. \url{https://lexicon.qumran-digital.org}}
\refapa{Rahlfs, A., \& Hanhart, R. (Eds.). (2006). \emph{Septuaginta} (Ed. rev.). Deutsche Bibelgesellschaft.}
\refapa{Rashi. (s.f.). \emph{Comentario a Éxodo y Proverbios}. Sefaria. \url{https://www.sefaria.org}}
\refapa{Reynolds, C.~W. (1987). Flocks, herds and schools: A distributed behavioral model. \emph{Computer Graphics, 21}(4), 25--34.}
\refapa{Sanders, J.~A. (1965). \emph{The Psalms scroll of Qumrân Cave 11 (11QPs\textsuperscript{a})} (Discoveries in the Judaean Desert IV). Clarendon Press.}
\refapa{\emph{Santa Biblia: Reina-Valera 1960}. (1960). Sociedades Bíblicas en América Latina.}
\refapa{Schreier, J. (Trad.). (2022). \emph{The Sefaria Midrash Rabbah} [Bereshit Rabá 1:1]. Sefaria. \url{https://www.sefaria.org}}
\refapa{Steinsaltz, A. (Ed.). (2017). \emph{Koren Talmud Bavli: Edición William Davidson} [Berajot 58a, 60b; Tamid 32a]. Sefaria. \url{https://www.sefaria.org}}
\refapa{Strogatz, S. (2003). \emph{Sync: The emerging science of spontaneous order}. Hyperion.}
\refapa{von Frisch, K. (1967). \emph{The dance language and orientation of bees}. Harvard University Press.}
\refapa{Windsor Clive, S., \& Smith, L. (2011). \emph{Murmuration} [Video]. Vimeo. \url{https://vimeo.com/31158841}}
\refapa{Wright, N.~T. (2013). \emph{Paul and the faithfulness of God}. Fortress Press.}

\end{document}
